% concatenated from main.tex
\documentclass[twoside,12pt]{article}

\usepackage[T1]{fontenc}
\usepackage[utf8]{inputenc}
\usepackage{lmodern}
\usepackage{microtype}
\usepackage{amsmath,amssymb,amsthm,mathtools}
\usepackage{xcolor}
\usepackage[a4paper,top=2cm,bottom=2cm,left=2cm,right=2cm]{geometry}
\usepackage{enumitem}
\usepackage[colorlinks=true,linkcolor=blue,citecolor=blue,urlcolor=blue]{hyperref}
\usepackage[nameinlink,noabbrev]{cleveref}

\hypersetup{
  pdftitle={Higher-Order Correlation Filters and Retained Spines for Multicolor Ramsey Numbers},
  pdfauthor={Yaping Mao and Gang Yang},
  pdfsubject={Upper bounds for multicolor diagonal Ramsey numbers},
  pdfkeywords={Ramsey numbers, book method, higher-order correlation, root filters, retained spines, relative entropy}
}

\allowdisplaybreaks
\emergencystretch=2em
\linespread{1.15}
\setlength{\parindent}{1.5em}
\setlength{\parskip}{0pt}
\setlist[enumerate]{itemsep=1pt,topsep=4pt,parsep=0pt,partopsep=0pt}

\newtheorem{theorem}{Theorem}[section]
\newtheorem{conjecture}[theorem]{Conjecture}
\newtheorem{lemma}[theorem]{Lemma}
\newtheorem{corollary}[theorem]{Corollary}
\newtheorem{claim}[theorem]{Claim}
\newtheorem{proposition}[theorem]{Proposition}

\theoremstyle{definition}
\newtheorem{definition}[theorem]{Definition}
\newtheorem{construction}[theorem]{Construction}

\theoremstyle{remark}
\newtheorem{remark}[theorem]{Remark}
\newtheorem{observation}[theorem]{Observation}

\crefname{theorem}{theorem}{theorems}
\Crefname{theorem}{Theorem}{Theorems}
\crefname{lemma}{lemma}{lemmas}
\Crefname{lemma}{Lemma}{Lemmas}
\crefname{proposition}{proposition}{propositions}
\Crefname{proposition}{Proposition}{Propositions}
\crefname{corollary}{corollary}{corollaries}
\Crefname{corollary}{Corollary}{Corollaries}
\crefname{definition}{definition}{definitions}
\Crefname{definition}{Definition}{Definitions}
\crefname{construction}{construction}{constructions}
\Crefname{construction}{Construction}{Constructions}
\crefname{claim}{claim}{claims}
\Crefname{claim}{Claim}{Claims}
\crefname{remark}{remark}{remarks}
\Crefname{remark}{Remark}{Remarks}
\crefname{observation}{observation}{observations}
\Crefname{observation}{Observation}{Observations}

\newcommand{\norm}[1]{\left\lVert #1\right\rVert}
\newcommand{\E}{\mathbb E}
\newcommand{\Pp}{\mathbb P}
\newcommand{\R}{\mathbb R}
\newcommand{\C}{\mathbb C}
\newcommand{\N}{\mathbb N}
\newcommand{\1}{\mathbf 1}
\newcommand{\ip}[2]{\left\langle #1,#2\right\rangle}
\newcommand{\tensor}{\otimes}
\newcommand{\dd}{\,\mathrm d}
\newcommand{\KL}{D}

\begin{document}

\title{\textbf{New upper bound
for multicolor Ramsey numbers}}
\author{Gang Yang\footnote{Graduate School of Environment and Information Sciences, Yokohama National University, 79-2 Tokiwadai, Hodogaya-ku, Yokohama 240-8501, Japan. {\tt gangyang98@outlook.com}},\quad Yaping Mao\footnote{Corresponding author. School of Mathematics and Statistics, Qinghai Normal University, Xining 810008, China. {\tt maoyaping@outlook.com}.}
}
\date{}
\maketitle

\begin{abstract}
Let $R_r(k)$ denote the diagonal $r$-color graph Ramsey number. We prove that
there exist absolute constants $c,K>0$ such that
\[
 R_r(k)\le
 \exp\!\left(-c\frac{k}{r^2\log^4(2r)}\right)r^{rk}
\]
for every $r\ge2$ and every $k\ge Kr^2\log^6(2r)$. The proof combines a
positive-coefficient root filter of variable order with a retained-spine
refinement of the multicolor book method.\\[2mm]
\noindent\textbf{Keywords:} Multicolor Ramsey numbers; book method;
higher-order correlation; root filters; relative entropy.

\medskip
\noindent\textbf{2020 Mathematics Subject Classification.}
05C55, 05D10, 30D15.
\end{abstract}
\section{Introduction}\label{sec:introduction}

For integers $r,k\ge2$, let $R_r(k)$ denote the least integer $n$ such that
every coloring of the edges of $K_n$ with $r$ colors contains a
monochromatic copy of $K_k$. Ramsey's theorem shows that $R_r(k)$ is finite
\cite{Ramsey1930}. The classical Erd\H{o}s--Szekeres recursion gives
$ R_r(k)\le r^{rk}$;
see \cite{ErdosSzekeres1935}. The purpose of this paper is to obtain an
exponential improvement whose dependence on $r$ remains effective when the
number of colors grows.

We use the following conventions throughout. All graphs are finite and simple,
all logarithms are natural, $\N=\{1,2,3,\ldots\}$, and
$[m]=\{1,\ldots,m\}$. For nonnegative quantities $f$ and $g$, the notation
$f=O(g)$ means that $f\le Cg$ for an absolute constant $C$, while
$f=\Omega(g)$ means that $f\ge cg$ for an absolute constant $c>0$. We write
$f=\Theta(g)$ when both estimates hold, that is, when there are absolute
constants $0<c\le C$ such that $cg\le f\le Cg$ throughout the stated range.
A subscript, as in $O_a(g)$, records the parameters on which the implicit
constant may depend. Finally, $f=o(g)$ in a specified limit means that
$f/g\to0$. We also write $(x)_+=\max\{x,0\}$, use $\1_A$ for the indicator
of a set or event $A$, and denote probability and expectation by $\Pp$ and
$\E$.

For two colors, Thomason, Conlon, and Sah obtained increasingly strong
subexponential improvements to the classical estimate
\cite{Conlon2009,Sah2023,Thomason1988}. Campos, Griffiths, Morris, and
Sahasrabudhe \cite{CamposEtAl2026} then proved the first exponential
improvement over the Erd\H{o}s--Szekeres upper bound. Gupta, Ndiaye, Norin,
and Wei \cite{GuptaEtAl2024} subsequently refined their argument and obtained
$R_2(k)\le3.8^{k+o(k)}$. We refer to
\cite{ConlonFoxSudakov2015,Wigderson2025} for broader accounts of graph
Ramsey numbers and the recent book method. Erd\H{o}s's probabilistic lower
bound \cite{Erdos1947} shows that the diagonal problem is exponential in
$k$, although a large gap between the upper and lower exponential rates
remains.

The multicolor case was subsequently treated by Balister, Bollob\'as,
Campos, Griffiths, Hurley, Morris, Sahasrabudhe, and Tiba
\cite{BalisterEtAl2026}. They proved that an exponential improvement holds
for every fixed number of colors and obtained the explicit estimate
\[
 R_r(k)\le \exp\!\left(-2^{-160}r^{-12}k\right)r^{rk}
 \qquad\bigl(k\ge 2^{200}r^{20}\bigr).
\]
In a recent preprint, Narang and Tang \cite{NarangTang2026} used robust OR
polynomials to sharpen the color dependence in the square-root correlation
estimate. Their result is
\begin{equation}\label{eq:narang-tang}
 R_r(k)\le
 \exp\!\left(-c\frac{k}{r^9\log^6(2r)}\right)r^{rk}
 \qquad\bigl(k\ge Cr^{14}\log^{12}(2r)\bigr)
\end{equation}
for absolute constants $c,C>0$.

Our main theorem gives a stronger dependence on $r$ in both the exponential
improvement and the range of $k$.

\begin{theorem}\label{thm:main}
There are absolute constants $c,K>0$ such that, for every $r\ge2$ and every
$k\ge Kr^2\log^6(2r)$, one has
\[
 R_r(k)\le
 \exp\!\left(-c\frac{k}{r^2\log^4(2r)}\right)r^{rk}.
\]
\end{theorem}

It is useful to compare the three bounds on the same scale. Balister et al.
obtained an exponential saving of order $k/r^{12}$ once $k$ is at least of
order $r^{20}$. Narang and Tang improved these two quantities to
$k/(r^9\log^6(2r))$ and $r^{14}\log^{12}(2r)$, respectively. The present
bound gives a saving of order $k/(r^2\log^4(2r))$ under the substantially
smaller condition $k\ge Kr^2\log^6(2r)$. Thus, relative to
\eqref{eq:narang-tang}, the saving in the exponent is larger by a factor of
order $r^7\log^2(2r)$, while the displayed sufficient lower bound on $k$ is
smaller by a factor of order $r^{12}\log^6(2r)$. These comparisons concern
the dependence on a growing number of colors. We do not try to optimize the
numerical base for a fixed small value of $r$; in particular, the specialized
two-color estimate of Gupta, Ndiaye, Norin, and Wei is stronger when $r=2$.

The proof yields the following more flexible statement. It is useful both for
tracking the dependence on the order of the root filter and for explaining the
choice of parameters in \Cref{thm:main}.

\begin{theorem}\label{thm:parameterised-main}
There are absolute constants $c,K>0$ such that the following holds. Let
$r\ge2$, and let $d$ be an integer satisfying
$
 3\le d\le\max\{3,\log(2r)\}.
$
If
\begin{equation}\label{eq:parameter-threshold}
 k\ge
 K r^{(2d^2-1)/(d-1)^2}
 d^{2d(2d-1)/(d-1)^2}
 \bigl(\log(2rd)\bigr)^{2d^2/(d-1)^2},
\end{equation}
then
\begin{equation}\label{eq:parameter-saving}
 R_r(k)\le
 \exp\!\left[
 -c\frac{k}{
 r^{2d/(d-1)}d^{4d/(d-1)}
 (\log(2rd))^{4/(d-1)}}
 \right]r^{rk}.
\end{equation}
\end{theorem}

For instance, $d=3$ gives an exponential improvement of order
$k/(r^3\log^2(2r))$, while $d=4$ gives one of order
$k/(r^{8/3}\log^{4/3}(2r))$. The choice
\[
 d=\max\left\{3,\left\lceil\frac12\log(2r)\right\rceil\right\}
\]
is admissible for every $r\ge2$ and satisfies $d=\Theta(\log(2r))$ as
$r\to\infty$. This choice gives \Cref{thm:main}. All constants in the
argument are independent of $d$ in the stated range.

The proof has two main new ingredients. The first is a higher-order
correlation lemma. Roughly speaking, it converts a failure of simultaneous
weak correlation among several vector-valued maps into a clustering
conclusion for one of the maps. In the correlation estimate used in the
earlier multicolor book argument, the relevant tail has a square-root
exponent. Our higher-order construction replaces this exponent by $1/d$,
where the integer $d$ may grow with the number of colors. This flexibility is
the source of the improved dependence on $r$.

The second ingredient is a refinement of the book argument. Recall that a
color-$i$ book consists of a color-$i$ clique, called its spine, together
with a set of common color-$i$ neighbors, called its page set. The initial
Erd\H{o}s--Szekeres regularization produces, for every color $i$, a
monochromatic clique $S_i$ joined in color $i$ to a common reservoir. Unlike
in the usual application of the book method, we retain all of these
preliminary cliques during the subsequent construction.

Suppose that the final book has color $i$, a spine of size $t$, and page set
$P$, and write $s_j=|S_j|$. The Ramsey problem remaining inside $P$ then has
target sizes
$
 k-s_1,\ldots,k-s_i-t,\ldots,k-s_r.
$
If $s_1+\cdots+s_r$ is large, the gain produced by regularization already
suffices. If it is small, then the distinguished target $k-s_i-t$ lies below
the diagonal by an amount of order $t$, and a one-coordinate entropy
estimate gives an additional factor
$\exp\!\left(-\Omega(t^2/k)\right)$. Combining this retained-spine gain with
the higher-order correlation lemma yields the stated bound.

The rest of the paper is organized as follows. In \Cref{sec:prelim} we
introduce the notation and collect the preliminary results used later. In
\Cref{sec:analytic} we prove the main analytic estimate. In \Cref{sec:book}
we use this estimate to establish the main combinatorial result. Finally, in
\Cref{sec:transfer} we combine these ingredients to prove the stated bounds
on $R_r(k)$.

\section{Preliminaries}\label{sec:prelim}

Given an $r$-edge-colored complete graph, a color $i\in[r]$, and a vertex
$v$, let
$
 N_i(v)=\{w:vw\text{ has color }i\}
$
be the color-$i$ neighborhood of $v$. If $X$ and $Y$ are nonempty vertex
sets, define the minimum color-$i$ density from $X$ to $Y$ by
\[
 p_i(X,Y)=\min_{x\in X}\frac{|N_i(x)\cap Y|}{|Y|}.
\]
The sets $X$ and $Y$ are not required to be disjoint.

A \emph{color-$i$ book} is an ordered pair $(T,P)$ of disjoint vertex sets
such that $T$ is a color-$i$ clique and every edge between $T$ and $P$ has
color $i$. We call $T$ the \emph{spine} and $P$ the \emph{page set}. A
$(t,m)$-book is a book satisfying $|T|=t$ and $|P|\ge m$. This terminology is
consistent with that used in the recent diagonal Ramsey work
\cite{CamposEtAl2026,Wigderson2025} and in its multicolor development
\cite{BalisterEtAl2026}.

For positive integers $k_1,\ldots,k_r$, let $R(k_1,\ldots,k_r)$ denote the
least $n$ with the following property: every $r$-edge-coloring of $K_n$
contains a color-$i$ copy of $K_{k_i}$ for some $i\in[r]$. For nonnegative
integers $b_1,\ldots,b_r$ whose sum is $B$, write
\[
 \binom{B}{b_1,\ldots,b_r}=\frac{B!}{b_1!\cdots b_r!}.
\]
Writing $B=k_1+\cdots+k_r$, the Erd\H{o}s--Szekeres recursion gives
\begin{equation}\label{eq:off-diagonal-standard}
 R(k_1,\ldots,k_r)
 \le \binom{B-r}{k_1-1,\ldots,k_r-1}.
\end{equation}
For later use we record the convenient weaker consequence
\begin{equation}\label{eq:off-diagonal-basic}
 R(k_1,\ldots,k_r)
 \le \binom{B}{k_1,\ldots,k_r}
 \le r^B.
\end{equation}
Indeed,
\[
 \frac{\binom{B}{k_1,\ldots,k_r}}
 {\binom{B-r}{k_1-1,\ldots,k_r-1}}
 =\frac{B(B-1)\cdots(B-r+1)}{\prod_{i=1}^r k_i}.
\]
Moreover,
\[
 B(B-1)\cdots(B-r+1)
 \ge(B-r+1)^r
 \ge\left(\frac Br\right)^r
 \ge\prod_{i=1}^r k_i,
\]
where the middle inequality follows from $B\ge r$ and the last one from the
arithmetic--geometric mean inequality. This proves the first inequality in
\eqref{eq:off-diagonal-basic}; the second follows from the multinomial
theorem.

We shall use the following form of the Erd\H{o}s--Szekeres regularization
lemma; see \cite[Lemma~5.2]{BalisterEtAl2026}. We refer to the sets $S_i$ as
the \emph{preliminary color-$i$ spines} and to $W$ as the \emph{common
reservoir}.

\begin{lemma}[\cite{BalisterEtAl2026}]\label{lem:regularization}
Let $\eta>0$, and let the edges of $K_n$ be colored with $r$ colors. There
are pairwise disjoint sets $S_1,\ldots,S_r,W$ such that, writing
$s_i=|S_i|$ and $s=\sum_i s_i$,
\begin{equation}\label{eq:regularization-size}
 |W|\ge\left(\frac{1+\eta}{r}\right)^s n,
\end{equation}
and, for every $w\in W$ and every $i\in[r]$,
\begin{equation}\label{eq:regularization-degree}
 |N_i(w)\cap W|
 \ge\left(\frac1r-\eta\right)|W|-1.
\end{equation}
Moreover, $S_i$ is a color-$i$ clique and every edge between $S_i$ and $W$
has color $i$.
\end{lemma}

For a real Hilbert space $\mathcal H$, we write
$\ip{x}{y}_{\mathcal H}$ for the inner product of $x,y\in\mathcal H$, and
$\norm{x}_{\mathcal H}=\sqrt{\ip{x}{x}_{\mathcal H}}$
for the induced norm. Whenever the underlying Hilbert space is clear from
the context, we omit the subscript.

Let $m\ge1$, and let $\mathcal K_1,\ldots,\mathcal K_m$ be real Hilbert
spaces. We write $\bigotimes_{j=1}^m\mathcal K_j$ for their Hilbert tensor
product over $\R$. If $x_j\in\mathcal K_j$ for every $j\in[m]$, then
$\bigotimes_{j=1}^m x_j=x_1\tensor\cdots\tensor x_m$
denotes the corresponding pure tensor. Its canonical inner product is
determined on pure tensors by
\[
 \ip{\bigotimes_{j=1}^m x_j}{\bigotimes_{j=1}^m y_j}_{\bigotimes_{j=1}^m\mathcal K_j}
 =\prod_{j=1}^m\ip{x_j}{y_j}_{\mathcal K_j}.
\]

For a nonnegative integer $a$ and a real Hilbert space $\mathcal H$, we write
\[
 \mathcal H^{\tensor a}=\bigotimes_{h=1}^a\mathcal H,
 \qquad
 x^{\tensor a}=\bigotimes_{h=1}^a x.
\]
For $a=0$, we use the conventions
$
 \mathcal H^{\tensor0}=\R$ and $
 x^{\tensor0}=1\in\R,
$
where $\R$ is equipped with its usual inner product. An empty product is
interpreted as $1$. In particular, if an exponent $a_i$ below is zero, then
the corresponding tensor factor and scalar factor are both interpreted as
$1$.

The following elementary positivity observation is the analytic starting
point.

\begin{lemma}\label{lem:positive-moments}
Let $r\ge1$ be an integer. Let $U$ and $U'$ be independent copies of a
random variable taking values in a finite set $\mathcal X$. For each
$i\in[r]$, let $\mathcal H_i$ be a real Hilbert space and let
$\sigma_i:\mathcal X\to\mathcal H_i$ be a map. Then, for all nonnegative
integers $a_1,\ldots,a_r$,
\begin{equation}\label{eq:positive-moments}
 \E\left(\prod_{i=1}^r
 \ip{\sigma_i(U)}{\sigma_i(U')}_{\mathcal H_i}^{a_i}\right)
 \ge0.
\end{equation}
\end{lemma}

\begin{proof}
Set
$
 \mathcal K=\bigotimes_{i=1}^r\mathcal H_i^{\tensor a_i}.
$
For every $x\in\mathcal X$, define
\[
Z(x)=\bigotimes_{i=1}^r\sigma_i(x)^{\tensor a_i}\in\mathcal K.
\]
For $x,y\in\mathcal X$, the defining property of the tensor-product inner
product gives
\begin{align*}
 \ip{Z(x)}{Z(y)}_{\mathcal K}
 =\prod_{i=1}^r
 \ip{\sigma_i(x)^{\tensor a_i}}{\sigma_i(y)^{\tensor a_i}}_{\mathcal H_i^{\tensor a_i}}=\prod_{i=1}^r
 \ip{\sigma_i(x)}{\sigma_i(y)}_{\mathcal H_i}^{a_i},
\end{align*}
where the identity remains valid when some $a_i=0$ by the conventions above.

Let $\mu(x)=\Pp(U=x)=\Pp(U'=x)$ for $x\in\mathcal X$. Since $U$ and $U'$
are independent copies,
\[
 \Pp(U=x,U'=y)=\mu(x)\mu(y)
 \qquad(x,y\in\mathcal X).
\]
Because $\mathcal X$ is finite, we may write the expectation as a finite sum:
\begin{align*}
 \E\left(\prod_{i=1}^r
 \ip{\sigma_i(U)}{\sigma_i(U')}_{\mathcal H_i}^{a_i}\right)
 &=\sum_{x,y\in\mathcal X}\mu(x)\mu(y)\ip{Z(x)}{Z(y)}_{\mathcal K}\\
 &=\ip{\sum_{x\in\mathcal X}\mu(x)Z(x)}
 {\sum_{y\in\mathcal X}\mu(y)Z(y)}_{\mathcal K}\\
 &=\norm{\sum_{x\in\mathcal X}\mu(x)Z(x)}_{\mathcal K}^{2}\ge0.
\end{align*}
This proves \eqref{eq:positive-moments}.
\end{proof}

\section{Higher-order correlation estimates}\label{sec:analytic}

Throughout this section, $\mathrm i$ denotes the imaginary unit, so that
$\mathrm i^2=-1$. For $z\in\C$, we write $\operatorname{Re}(z)$ for the
real part of $z$.

\subsection{The root filter}\label{sec:filters}

Recall that an \emph{entire function} is a complex-valued function that is
holomorphic on all of $\C$. We call an entire function
\emph{positive-coefficient} if every coefficient in its Maclaurin expansion
is a nonnegative real number. Fix an integer $d\ge3$, and define
\[
 E_d(z)=\sum_{n=0}^{\infty}\frac{z^n}{(dn)!}.
\]
The factorial denominator gives the series infinite radius of convergence.
Thus $E_d$ is an entire function with nonnegative Taylor coefficients. Set
\[
 \theta_d=\frac{\pi}{d},
 \qquad
 \eta_d=1-\cos(2\theta_d),
 \qquad
 \gamma_d=1-\cos\theta_d.
\]

\begin{lemma}\label{lem:root-unity}
Let $d\ge3$ be an integer, and let
\[
 \omega_d=\exp\!\left(\frac{2\pi\mathrm i}{d}\right).
\]
Then, for every $u\ge0$,
\begin{align}
 E_d(u^d)
 &=\frac1d\sum_{j=0}^{d-1}\exp(u\omega_d^j),
 \label{eq:positive-root-filter}\\
 E_d(-u^d)
 &=\frac1d\sum_{j=0}^{d-1}
 \exp\!\left(u\exp\!\left(\frac{(2j+1)\pi\mathrm i}{d}\right)\right).
 \label{eq:negative-root-filter}
\end{align}
Consequently, for every $u\ge0$,
\begin{equation}\label{eq:negative-axis-bound}
 |E_d(-u^d)|\le e^{u\cos\theta_d}
\end{equation}
and
\begin{equation}\label{eq:positive-axis-upper}
 E_d(u^d)\le e^u.
\end{equation}
Moreover, for every $u\ge0$ satisfying
\begin{equation}\label{eq:positive-axis-condition}
 u\eta_d\ge\log\bigl(2(d-1)\bigr),
\end{equation}
one has
\begin{equation}\label{eq:positive-axis-lower}
 E_d(u^d)\ge\frac{e^u}{2d}.
\end{equation}
\end{lemma}

\begin{proof}
For every nonnegative integer $m$, the standard roots-of-unity identity gives
\begin{equation}\label{eq:root-orthogonality}
 \frac1d\sum_{j=0}^{d-1}\omega_d^{jm}
 =
 \begin{cases}
  1,&d\mid m,\\
  0,&d\nmid m.
 \end{cases}
\end{equation}
Indeed, if $d\mid m$, every summand is $1$. If $d\nmid m$, then
$\omega_d^m\ne1$, and the geometric-sum formula gives
\[
 \sum_{j=0}^{d-1}\omega_d^{jm}
 =\frac{1-\omega_d^{dm}}{1-\omega_d^m}=0.
\]
Since the exponential series converges absolutely and the sum over $j$ is
finite, we may interchange the two sums. Using
\eqref{eq:root-orthogonality}, we obtain
\begin{align*}
 \frac1d\sum_{j=0}^{d-1}\exp(u\omega_d^j)
=\sum_{m=0}^{\infty}\frac{u^m}{m!}
 \left(\frac1d\sum_{j=0}^{d-1}\omega_d^{jm}\right)=\sum_{n=0}^{\infty}\frac{u^{dn}}{(dn)!}=E_d(u^d).
\end{align*}
This proves \eqref{eq:positive-root-filter}.

For $j=0,\ldots,d-1$, put
\[
 \zeta_j
 =\exp\!\left(\frac{(2j+1)\pi\mathrm i}{d}\right)
 =\exp\!\left(\frac{\pi\mathrm i}{d}\right)\omega_d^j.
\]
The numbers $\zeta_0,\ldots,\zeta_{d-1}$ are precisely the $d$th roots of
$-1$. For every nonnegative integer $m$,
\[
 \frac1d\sum_{j=0}^{d-1}\zeta_j^m
 =\exp\!\left(\frac{\pi\mathrm i m}{d}\right)
 \left(\frac1d\sum_{j=0}^{d-1}\omega_d^{jm}\right).
\]
By \eqref{eq:root-orthogonality}, this expression vanishes unless $m=dn$
for some nonnegative integer $n$. If $m=dn$, it equals
$\exp(\pi\mathrm i n)=(-1)^n$. Therefore
\begin{align*}
 \frac1d\sum_{j=0}^{d-1}\exp(u\zeta_j)
=\sum_{m=0}^{\infty}\frac{u^m}{m!}
 \left(\frac1d\sum_{j=0}^{d-1}\zeta_j^m\right)=\sum_{n=0}^{\infty}\frac{(-1)^nu^{dn}}{(dn)!}=E_d(-u^d),
\end{align*}
which proves \eqref{eq:negative-root-filter}.

For every $j\in\{0,\ldots,d-1\}$,
\[
 \operatorname{Re}(\zeta_j)
 =\cos\!\left(\frac{(2j+1)\pi}{d}\right)
 \le\cos\!\left(\frac\pi d\right)=\cos\theta_d.
\]
Hence,
$
 |\exp(u\zeta_j)|
=\exp\!\left(u\operatorname{Re}(\zeta_j)\right)
 \le e^{u\cos\theta_d}.
$
Applying the triangle inequality to \eqref{eq:negative-root-filter} proves
\eqref{eq:negative-axis-bound}.

It remains to prove the positive-axis estimates. Since all terms in the
defining series of $E_d(u^d)$ are nonnegative,
\[
 E_d(u^d)
 =\sum_{n=0}^{\infty}\frac{u^{dn}}{(dn)!}
 \le\sum_{m=0}^{\infty}\frac{u^m}{m!}=e^u,
\]
which proves \eqref{eq:positive-axis-upper}. For the lower bound, in
\eqref{eq:positive-root-filter} the term corresponding to $j=0$ is $e^u$.
For every $j\in\{1,\ldots,d-1\}$,
\[
 \operatorname{Re}(\omega_d^j)
 =\cos\!\left(\frac{2\pi j}{d}\right)
 \le\cos\!\left(\frac{2\pi}{d}\right)=\cos(2\theta_d).
\]
Since $E_d(u^d)$ is real, taking real parts in
\eqref{eq:positive-root-filter} and using the triangle inequality gives
\[
 dE_d(u^d)
 \ge e^u-(d-1)e^{u\cos(2\theta_d)}
 =e^u\left(1-(d-1)e^{-u\eta_d}\right).
\]
Condition \eqref{eq:positive-axis-condition} implies
$(d-1)e^{-u\eta_d}\le1/2$, and hence
$E_d(u^d)\ge e^u/(2d)$. This proves
\eqref{eq:positive-axis-lower}.
\end{proof}

Define
\[
 u_{r,d}=\max\left\{
 1,
 \frac{\log(2(d-1))}{\eta_d},
 \frac{\log(24d^2r)}{2\gamma_d}
 \right\},
 \qquad
 L_{r,d}=u_{r,d}^d,
\]
and
$a_{r,d}=\exp(-2u_{r,d}\cos\theta_d)$.
The inequalities
\[
 1-\cos x\ge\frac{2x^2}{\pi^2}
 \qquad(0\le x\le\pi)
\]
and
\[
 1-\cos(2x)\ge\frac{8x^2}{\pi^2}
 \qquad(0\le x\le\pi/2)
\]
give
\[
 \gamma_d\ge\frac{2}{d^2},
 \qquad
 \eta_d\ge\frac{8}{d^2}.
\]
Consequently, there is an absolute constant $C_0>0$ such that
\begin{equation}\label{eq:urd-bound}
 u_{r,d}\le C_0d^2\log(2rd).
\end{equation}

Define
\[
 H_{r,d}(z)=1+a_{r,d}E_d(z)^2,
 \qquad
 G_{r,d}(z)=\frac{H_{r,d}(z)-H_{r,d}(-z)}2.
\]
The Taylor coefficients of $H_{r,d}$ are nonnegative, and $G_{r,d}$ is the
odd part of $H_{r,d}$, so its Taylor coefficients are nonnegative as well.
Also, $H_{r,d}(x)\ge1$ for every real $x$.

\begin{lemma}\label{lem:ratio-estimates}
Let $r\ge2$ and $d\ge3$ be integers. Then, for every $x\ge0$,
\[
 0\le\frac{G_{r,d}(x)}{H_{r,d}(x)}\le\frac12.
\]
For every $x<0$,
\[
 \frac{G_{r,d}(x)}{H_{r,d}(x)}\le0.
\]
Moreover, if $y\ge L_{r,d}$, then
\[
 \frac{G_{r,d}(-y)}{H_{r,d}(-y)}\le-r.
\]
\end{lemma}

\begin{proof}
Since $E_d(x)$ is real for every real $x$,
\[
 H_{r,d}(x)=1+a_{r,d}E_d(x)^2\ge1.
\]
Let $x\ge0$. The nonnegativity of the Taylor coefficients of $E_d$ gives
$|E_d(-x)|\le E_d(x)$.
Therefore, $1\le H_{r,d}(-x)\le H_{r,d}(x)$,
and hence
\[
 \frac{G_{r,d}(x)}{H_{r,d}(x)}
 =\frac12\left(1-\frac{H_{r,d}(-x)}{H_{r,d}(x)}\right)
 \in\left[0,\frac12\right].
\]

If $x<0$, write $x=-z$ with $z>0$. Since $G_{r,d}$ is odd and
$G_{r,d}(z)\ge0$, we have $G_{r,d}(x)\le0$. Because $H_{r,d}(x)>0$, the
second assertion follows.

Now let $y\ge L_{r,d}$. Write $y=u^d$, where $u\ge u_{r,d}$. By the
definition of $u_{r,d}$ and \Cref{lem:root-unity},
\[
 E_d(y)\ge\frac{e^u}{2d}.
\]
Thus
\begin{equation}\label{eq:ratio-H-positive}
 H_{r,d}(y)
 \ge a_{r,d}E_d(y)^2
 \ge a_{r,d}\frac{e^{2u}}{4d^2}.
\end{equation}
On the other hand, \eqref{eq:negative-axis-bound} gives
\[
 H_{r,d}(-y)
 \le1+a_{r,d}e^{2u\cos\theta_d}.
\]
Since $u\ge u_{r,d}$,
$a_{r,d}e^{2u\cos\theta_d}
 =e^{2(u-u_{r,d})\cos\theta_d}\ge1$,
and therefore
\begin{equation}\label{eq:ratio-H-negative}
 H_{r,d}(-y)\le2a_{r,d}e^{2u\cos\theta_d}.
\end{equation}
Combining \eqref{eq:ratio-H-positive} and
\eqref{eq:ratio-H-negative}, we obtain
\[
 \frac{H_{r,d}(y)}{H_{r,d}(-y)}
 \ge\frac{e^{2u(1-\cos\theta_d)}}{8d^2}
 =\frac{e^{2\gamma_du}}{8d^2}.
\]
By the definition of $u_{r,d}$, the last expression is at least $3r$.
Finally,
\[
 \frac{G_{r,d}(-y)}{H_{r,d}(-y)}
 =-\frac12\left(\frac{H_{r,d}(y)}{H_{r,d}(-y)}-1\right)
 \le-\frac{3r-1}{2}\le-r.
\]
This completes the proof.
\end{proof}

We now combine these one-variable functions in $r$ coordinates. Set
\[
 F_{r,d}(x_1,\ldots,x_r)
 =\sum_{j=1}^rG_{r,d}(x_j)\prod_{i\ne j}H_{r,d}(x_i).
\]
Every multivariate Taylor coefficient of $F_{r,d}$ is nonnegative, and
\begin{equation}\label{eq:separator-factorisation}
 F_{r,d}(x_1,\ldots,x_r)
 =\left(\prod_{i=1}^rH_{r,d}(x_i)\right)
 \left(\sum_{j=1}^r\frac{G_{r,d}(x_j)}{H_{r,d}(x_j)}\right).
\end{equation}

\begin{lemma}\label{lem:separator-bounds}
Let $r\ge2$ and $d\ge3$ be integers. The following statements hold.
\begin{enumerate}[label=\textup{(\roman*)}]
\item Let $x=(x_1,\ldots,x_r)\in\R^r$. If $x_k\le-L_{r,d}$ for some
$k\in[r]$, then
$F_{r,d}(x_1,\ldots,x_r)\le-1.$

\item Let $A_1,\ldots,A_r\ge-1$, and put
$M=\max_{i\in[r]}A_i$. Then
\begin{equation}\label{eq:separator-growth}
 F_{r,d}(L_{r,d}A_1,\ldots,L_{r,d}A_r)
 \le D_r\exp\!\left(c_{r,d}(M+1)^{1/d}\right),
\end{equation}
where
$
 D_r=r2^{r-1}$ and $
 c_{r,d}=2ru_{r,d}.
$
\end{enumerate}
\end{lemma}

\begin{proof}
For every real $z$,
\begin{equation}\label{eq:H-at-least-one}
 H_{r,d}(z)\ge1.
\end{equation}
Suppose first that $x_k\le-L_{r,d}$ for some $k\in[r]$. By
\Cref{lem:ratio-estimates},
\[
 \frac{G_{r,d}(x_k)}{H_{r,d}(x_k)}\le-r.
\]
For every $j\ne k$, the same lemma gives
\[
 \frac{G_{r,d}(x_j)}{H_{r,d}(x_j)}\le\frac12.
\]
Consequently,
\[
 \sum_{j=1}^r\frac{G_{r,d}(x_j)}{H_{r,d}(x_j)}
 \le-r+\frac{r-1}{2}=-\frac{r+1}{2}\le-1.
\]
Using \eqref{eq:separator-factorisation} and
\eqref{eq:H-at-least-one}, we obtain
$
 F_{r,d}(x_1,\ldots,x_r)\le-1.
$
This proves part~\textup{(i)}.

For part~\textup{(ii)}, we first establish the one-variable estimate
\begin{equation}\label{eq:H-growth}
 H_{r,d}(L_{r,d}A)
 \le2\exp\!\left(2u_{r,d}(A+1)^{1/d}\right)
 \qquad(A\ge-1).
\end{equation}
Suppose first that $-1\le A\le0$. Define
$u=u_{r,d}(-A)^{1/d}$. Then $0\le u\le u_{r,d}$ and
$L_{r,d}A=-u^d$. By \eqref{eq:negative-axis-bound},
$
 a_{r,d}|E_d(-u^d)|^2
 \le e^{-2u_{r,d}\cos\theta_d}e^{2u\cos\theta_d}\le1.
$
Thus $H_{r,d}(L_{r,d}A)\le2$, which implies
\eqref{eq:H-growth} in this range.

Now suppose that $A\ge0$. Define $u=u_{r,d}A^{1/d}$, so that
$L_{r,d}A=u^d$. Equation \eqref{eq:positive-axis-upper} gives
$
 E_d(L_{r,d}A)\le e^{u_{r,d}A^{1/d}}.
$
Since $0<a_{r,d}\le1$,
\[
 H_{r,d}(L_{r,d}A)
 \le1+e^{2u_{r,d}A^{1/d}}
 \le2e^{2u_{r,d}(A+1)^{1/d}}.
\]
This proves \eqref{eq:H-growth} for all $A\ge-1$.

Put $\widetilde x_i=L_{r,d}A_i$. If
$F_{r,d}(\widetilde x_1,\ldots,\widetilde x_r)<0$, then
\eqref{eq:separator-growth} is immediate. We may therefore assume that this
quantity is nonnegative. By \eqref{eq:separator-factorisation},
\[
 0\le\sum_{i=1}^r
 \frac{G_{r,d}(\widetilde x_i)}{H_{r,d}(\widetilde x_i)}
 \le\frac r2.
\]
Using \eqref{eq:H-growth}, we obtain
\begin{align*}
 F_{r,d}(\widetilde x_1,\ldots,\widetilde x_r)
 \le\frac r2\prod_{i=1}^rH_{r,d}(\widetilde x_i)\le r2^{r-1}\exp\!\left(
2u_{r,d}\sum_{i=1}^r(A_i+1)^{1/d}\right)\le r2^{r-1}\exp\!\left(2ru_{r,d}(M+1)^{1/d}\right).
\end{align*}
This is \eqref{eq:separator-growth}.
\end{proof}

\subsection{A higher-order correlation theorem}\label{sec:correlation}

\begin{theorem}\label{thm:correlation}
Let $r\ge2$ and $d\ge3$ be integers. Let $U$ and $U'$ be independent
copies of a random variable taking values in a finite set $\mathcal X$.
For each $i\in[r]$, let $\mathcal H_i$ be a real Hilbert space and let
$\sigma_i:\mathcal X\to\mathcal H_i$ be an arbitrary map. Set
$
 Z_i=\ip{\sigma_i(U)}{\sigma_i(U')}.
$
Then there exist $i\in[r]$ and $\lambda\ge-1$ such that
\begin{equation}\label{eq:correlation-tail}
 \Pp\bigl(
  Z_i\ge\lambda,\
  Z_j\ge-1\text{ for every }j\ne i
 \bigr)
 \ge
 \beta_r\exp\!\left(-C_{r,d}(\lambda+1)^{1/d}\right),
\end{equation}
where
\begin{equation}\label{eq:correlation-constants}
 \beta_r=\frac1{4D_r(r+1)},
 \qquad
 C_{r,d}=2c_{r,d}=4ru_{r,d}.
\end{equation}
In particular, there are absolute constants $C_1,C_2>0$ such that
\begin{equation}\label{eq:correlation-constant-bounds}
 C_{r,d}\le C_1rd^2\log(2rd),
 \qquad
 \log(1/\beta_r)\le C_2r.
\end{equation}
\end{theorem}

\begin{proof}
For brevity, write
$
 L=L_{r,d},
 c=c_{r,d},$ and $
 C=C_{r,d}.
$
Since $\mathcal X$ is finite, each $Z_i$ has finite range. In particular,
the random vector $(LZ_1,\ldots,LZ_r)$ is bounded.

Write the multivariate Taylor expansion of $F_{r,d}$ as
\[
 F_{r,d}(x_1,\ldots,x_r)
 =\sum_{\mathbf a\in\mathbb Z_{\ge0}^{\,r}}
 q_{\mathbf a}\prod_{i=1}^r x_i^{a_i},
 \qquad q_{\mathbf a}\ge0,
\]
where $\mathbf a=(a_1,\ldots,a_r)$ and
$|\mathbf a|=a_1+\cdots+a_r$. To justify taking expectations term by term,
set
\[
 B=\max\left\{
 1,
 L\max_{\substack{i\in[r]\\x,y\in\mathcal X}}
 \left|\ip{\sigma_i(x)}{\sigma_i(y)}\right|
 \right\}.
\]
Then $|LZ_i|\le B$ almost surely for every $i\in[r]$. Since $F_{r,d}$ is
entire and its Taylor coefficients are nonnegative,
\[
 \sum_{\mathbf a\in\mathbb Z_{\ge0}^{\,r}}
 q_{\mathbf a}L^{|\mathbf a|}
 \E\left(\left|\prod_{i=1}^rZ_i^{a_i}\right|\right)
 \le
 \sum_{\mathbf a\in\mathbb Z_{\ge0}^{\,r}}
 q_{\mathbf a}B^{|\mathbf a|}
 =F_{r,d}(B,\ldots,B)<\infty.
\]
Thus the Taylor series is absolutely summable in $L^1$, so expectation may
be interchanged with summation. By \Cref{lem:positive-moments},
\begin{equation}\label{eq:separator-positive-expectation}
 \E\bigl(F_{r,d}(LZ_1,\ldots,LZ_r)\bigr)
 =\sum_{\mathbf a}q_{\mathbf a}L^{|\mathbf a|}
 \E\left(\prod_{i=1}^rZ_i^{a_i}\right)
 \ge0.
\end{equation}

Let
$\mathcal E=\{Z_j\ge-1\text{ for every }j\in[r]\}$, $\bar{\mathcal E}=\{Z_j<-1\text{ for some }j\in[r]\}$
and
 $M=\max_{j\in[r]}Z_j$,
and define
\[
 Y=\bigl((M+1)_+\bigr)^{1/d},
 \qquad
 (t)_+=\max\{t,0\}.
\]
Thus $Y\ge0$, and on $\mathcal E$ we have $Y=(M+1)^{1/d}$.

If $\bar{\mathcal E}$ occurs, then $Z_k<-1$ for some $k\in[r]$, and hence
$LZ_k<-L$. Therefore, \Cref{lem:separator-bounds}(i) gives
\[
 F_{r,d}(LZ_1,\ldots,LZ_r)\le-1
 \qquad\text{on }\bar{\mathcal E}.
\]
On the other hand, on $\mathcal E$ all coordinates $Z_j$ are at least
$-1$. Applying \Cref{lem:separator-bounds}(ii) with $A_j=Z_j$ gives
\[
 F_{r,d}(LZ_1,\ldots,LZ_r)\le D_re^{cY}
 \qquad\text{on }\mathcal E.
\]
Consequently,
\[
 F_{r,d}(LZ_1,\ldots,LZ_r)
 \le D_re^{cY}\1_{\mathcal E}-\1_{\bar{\mathcal E}}.
\]
Taking expectations and using
\eqref{eq:separator-positive-expectation}, we obtain
\begin{equation}\label{eq:correlation-master}
 D_r\E\bigl[e^{cY}\1_{\mathcal E}\bigr]
 \ge1-\Pp(\mathcal E).
\end{equation}

For $i\in[r]$ and $t\ge-1$, put
$
 \mathcal A_i(t)
 =\{Z_i\ge t,\ Z_j\ge-1\text{ for every }j\ne i\}.
$
If $\Pp(\mathcal E)\ge\beta_r$, choose $\lambda=-1$. Then
$\mathcal A_i(-1)=\mathcal E$ for every $i\in[r]$, while
$(\lambda+1)^{1/d}=0$. Hence \eqref{eq:correlation-tail} follows.

It remains to consider the case $\Pp(\mathcal E)<\beta_r$. Suppose, for a
contradiction, that \eqref{eq:correlation-tail} fails for every $i\in[r]$
and every $\lambda\ge-1$. Thus
\begin{equation}\label{eq:correlation-failure}
 \Pp\bigl(\mathcal A_i(\lambda)\bigr)
 <\beta_re^{-C(\lambda+1)^{1/d}}
\end{equation}
for every $i\in[r]$ and every $\lambda\ge-1$.

Fix $u\ge0$ and set $\lambda_u=u^d-1$. Then $\lambda_u\ge-1$ and
$(\lambda_u+1)^{1/d}=u$. Moreover, on $\mathcal E$,
 $Y\ge u$ if and only if
 $M\ge\lambda_u$.
Therefore
\[
 \mathcal E\cap\{Y\ge u\}
 =\bigcup_{i=1}^r\mathcal A_i(\lambda_u).
\]
The union bound and \eqref{eq:correlation-failure} yield
\begin{equation}\label{eq:correlation-Y-tail}
 \Pp\bigl(\mathcal E\cap\{Y\ge u\}\bigr)
 <r\beta_re^{-Cu}.
\end{equation}

For every $y\ge0$,
$e^{cy}=1+c\int_0^\infty e^{cu}\1_{\{y>u\}}\,\dd u.$
Multiplying by $\1_{\mathcal E}$, taking expectations, and applying
Tonelli's theorem gives
\begin{align*}
 \E\bigl[e^{cY}\1_{\mathcal E}\bigr]
 &=\Pp(\mathcal E)
 +c\int_0^\infty e^{cu}
 \Pp\bigl(\mathcal E\cap\{Y>u\}\bigr)\,\dd u\\
 &\le\Pp(\mathcal E)
 +c\int_0^\infty e^{cu}
 \Pp\bigl(\mathcal E\cap\{Y\ge u\}\bigr)\,\dd u\\
 &<\beta_r+r\beta_rc
 \int_0^\infty e^{-(C-c)u}\,\dd u\\
 &=\beta_r+r\beta_r\frac{c}{C-c}
 =(r+1)\beta_r
 =\frac1{4D_r}.
\end{align*}
Here we used $C=2c>c>0$. Consequently,
$D_r\E\bigl[e^{cY}\1_{\mathcal E}\bigr]<\frac14$.
On the other hand, \eqref{eq:correlation-master} and
$\Pp(\mathcal E)<\beta_r$ give
$D_r\E\bigl[e^{cY}\1_{\mathcal E}\bigr]
 >1-\beta_r$.
Since $r\ge2$ and $D_r=r2^{r-1}\ge4$,
\[
 \beta_r=\frac1{4D_r(r+1)}\le\frac1{48}<\frac14,
\]
and hence $1-\beta_r>3/4$, a contradiction. This proves
\eqref{eq:correlation-tail}.

Finally, \eqref{eq:urd-bound} and $C_{r,d}=4ru_{r,d}$ give
$C_{r,d}\le4C_0rd^2\log(2rd),$
so one may take $C_1=4C_0$. Moreover,
\[
 \log(1/\beta_r)
 =\log\bigl(4r2^{r-1}(r+1)\bigr)
 =(r+1)\log2+\log r+\log(r+1)
 \le3r
\]
for every $r\ge2$, so one may take $C_2=3$.
\end{proof}

\section{Density increments and monochromatic books}\label{sec:book}

\subsection{A density-increment lemma}\label{sec:density}

For later use, we record the correlation theorem in a parameterized form.
Let $r\ge2$ and $d\ge3$ be integers, and let $0<\beta\le1$ and $C>0$.
We say that $\mathcal G_r^{(d)}(\beta,C)$ holds if the following property is
satisfied. For every finite set $\mathcal X$, every pair $U,U'$ of
independent copies of a random variable taking values in $\mathcal X$, and
every collection of maps $\sigma_i:\mathcal X\to\mathcal H_i$ into real
Hilbert spaces, set
$
 Z_i=\ip{\sigma_i(U)}{\sigma_i(U')}
$, where $i\in[r].
$
Then there exist $i\in[r]$ and $\lambda\ge-1$ such that
\[
 \Pp\bigl(
  Z_i\ge\lambda,
  Z_j\ge-1\text{ for every }j\ne i
 \bigr)
 \ge\beta\exp\!\left(-C(\lambda+1)^{1/d}\right).
\]
Thus \Cref{thm:correlation} states that
$\mathcal G_r^{(d)}(\beta_r,C_{r,d})$ holds.

Recall that, for nonempty finite vertex sets $A$ and $B$,
\[
 p_i(A,B)=\min_{v\in A}\frac{|N_i(v)\cap B|}{|B|}
\]
is the minimum relative color-$i$ degree from $A$ into $B$.

\begin{lemma}\label{lem:key}
Let $r\ge2$ and $d\ge3$ be integers, let $0<\beta\le1$ and $C>0$, and
assume that $\mathcal G_r^{(d)}(\beta,C)$ holds. Let
$X,Y_1,\ldots,Y_r$ be nonempty finite vertex sets in an $r$-edge-colored
complete graph, and put
$
 p_i=p_i(X,Y_i)>0
~(i\in[r]).
$
For every choice of positive real numbers $\alpha_1,\ldots,\alpha_r$, there
exist a vertex $x\in X$, an index $\ell\in[r]$, a real number
$\lambda\ge-1$, a nonempty set $X'\subseteq X$, and nonempty sets
$
 Y_i'\subseteq N_i(x)\cap Y_i
~(i\in[r])
$
such that
\begin{equation}\label{eq:key-X}
 |X'|\ge\beta e^{-C(\lambda+1)^{1/d}}|X|,
\end{equation}
\begin{equation}\label{eq:key-selected}
 p_\ell(X',Y_\ell')\ge p_\ell+\lambda\alpha_\ell,
\end{equation}
and, for every $i\in[r]$,
$
 |Y_i'|=p_i|Y_i|,
 p_i(X',Y_i')\ge p_i-\alpha_i.
$
\end{lemma}

\begin{proof}
For each $i\in[r]$, set
$
 m_i=p_i|Y_i|
 =\min_{v\in X}|N_i(v)\cap Y_i|.
$
Since $X$ and $Y_i$ are finite and $p_i>0$, the number $m_i$ is a positive
integer. For every $v\in X$, choose a set
$A_i(v)\subseteq N_i(v)\cap Y_i,
 |A_i(v)|=m_i.
$
Regard the indicator of a subset of $Y_i$ as a vector in the real Hilbert
space $\R^{Y_i}$ with its standard inner product. Define
\[
 \sigma_i(v)=
 \frac{\1_{A_i(v)}-p_i\1_{Y_i}}
 {\sqrt{\alpha_ip_i|Y_i|}},
 \qquad v\in X.
\]
The denominator is nonzero because $\alpha_i>0$, $p_i>0$, and
$Y_i\ne\emptyset$. For $u,v\in X$, using
$|A_i(u)|=|A_i(v)|=p_i|Y_i|$, we obtain
\begin{equation}\label{eq:key-inner-product}
 \ip{\sigma_i(u)}{\sigma_i(v)}
 =\frac{|A_i(u)\cap A_i(v)|-p_i^2|Y_i|}
 {\alpha_ip_i|Y_i|}.
\end{equation}

Let $U,U'$ be independent uniform random variables on $X$. By
$\mathcal G_r^{(d)}(\beta,C)$, there exist $\ell\in[r]$ and
$\lambda\ge-1$ such that, on writing
$q=\beta\exp\!\left(-C(\lambda+1)^{1/d}\right)$,
we have
\[
 \Pp\bigl(
  \ip{\sigma_\ell(U)}{\sigma_\ell(U')}\ge\lambda,
  \ip{\sigma_i(U)}{\sigma_i(U')}\ge-1
  \text{ for every }i\ne\ell
 \bigr)\ge q.
\]
For each $u\in X$, define
\[
 X_u=\left\{v\in X:
 \begin{array}{l}
  \ip{\sigma_\ell(u)}{\sigma_\ell(v)}\ge\lambda,
  \ip{\sigma_i(u)}{\sigma_i(v)}\ge-1
  \text{ for every }i\ne\ell
 \end{array}
 \right\}.
\]
Since $U$ and $U'$ are independent and uniform on $X$,
\[
 q\le\frac1{|X|}\sum_{u\in X}\frac{|X_u|}{|X|}.
\]
Consequently, there exists $x\in X$ such that $|X_x|\ge q|X|$. Set
$
 X'=X_x,
 Y_i'=A_i(x)
~(i\in[r]).
$
Then \eqref{eq:key-X} holds, and
$
 Y_i'\subseteq N_i(x)\cap Y_i,
 |Y_i'|=p_i|Y_i|.
$
In particular, $X'$ and all the sets $Y_i'$ are nonempty.

For $i\in[r]$, define
\[
 \tau_i=
 \begin{cases}
  \lambda,&i=\ell,\\
  -1,&i\ne\ell.
 \end{cases}
\]
For every $v\in X'$, the definition of $X'$ gives
$\ip{\sigma_i(x)}{\sigma_i(v)}\ge\tau_i$. Using
\eqref{eq:key-inner-product} and $Y_i'=A_i(x)$, we obtain
$
 |A_i(v)\cap Y_i'|
 \ge p_i^2|Y_i|+\tau_i\alpha_ip_i|Y_i|.
$
Since $A_i(v)\subseteq N_i(v)$, division by
$|Y_i'|=p_i|Y_i|$ gives
\[
 \frac{|N_i(v)\cap Y_i'|}{|Y_i'|}
 \ge p_i+\tau_i\alpha_i.
\]
Taking the minimum over $v\in X'$ yields
$
 p_i(X',Y_i')\ge p_i+\tau_i\alpha_i.
$
For $i=\ell$, this is \eqref{eq:key-selected}. Since $\lambda\ge-1$, we
also have $\tau_i\ge-1$ for every $i$, and hence
$p_i(X',Y_i')\ge p_i-\alpha_i$.
\end{proof}

\subsection{The book lemma}\label{sec:book-theorem}

\begin{theorem}\label{thm:book}
Let $r\ge2$ and $d\ge3$ be integers. Let $0<\beta\le1$ and $C>0$, and
assume that $\mathcal G_r^{(d)}(\beta,C)$ holds. Let $t,m$ be positive
integers, let $0<p\le1$, and let $\delta,\lambda_0>0$. Set
\[
 L=\log(1/\delta),
 \qquad
 L_p=\log(1/p),
 \qquad
 \rho=\log(r/\beta),
\]
and suppose that
\[
 0<\delta\le\min\{p/4,1/4\},
 \qquad
 \lambda_0\ge\max\{2,6L\},
 \qquad
 t\ge\lambda_0\delta^{-1/(d-1)}.
\]
Define
\[
 \Pi=\frac{3\delta}{p}+\frac{6LL_p}{\lambda_0},
 \qquad
 \Xi=2\rho+4C\lambda_0^{1/d}
      +\frac{12CL}{\lambda_0^{(d-1)/d}}.
\]
Let $X,Y_1,\ldots,Y_r$ be nonempty vertex sets in an $r$-edge-colored
complete graph. Suppose that
\begin{equation}\label{eq:book-density-hyp}
 |N_i(x)\cap Y_i|\ge p|Y_i|
 \qquad(x\in X,\ i\in[r]),
\end{equation}
\begin{equation}\label{eq:book-page-hyp}
 |Y_i|\ge p^{-t}e^{\Pi t}m
 \qquad(i\in[r]),
\end{equation}
and
\begin{equation}\label{eq:book-reservoir-hyp}
 |X|\ge2rt\,e^{rt\Xi}.
\end{equation}
Then there exist an index $i\in[r]$, a color-$i$ clique $T\subseteq X$
with $|T|=t$, and a set $P\subseteq Y_i\setminus T$ with $|P|=m$ such
that every edge between $T$ and $P$ has color $i$. In particular, $(T,P)$
is a color-$i$ $(t,m)$-book.
\end{theorem}

\begin{proof}
We construct the spines one vertex at a time. At time $s$, let $X(s)$ be the
current reservoir, let $Y_i(s)$ be the current page set for color $i$, and
let $T_i(s)$ be the color-$i$ spine constructed so far. Initially,
$
 X(0)=X,
 Y_i(0)=Y_i,$ and $
 T_i(0)=\emptyset.
$
Whenever the relevant sets are nonempty, write
$p_i(s)=p_i\bigl(X(s),Y_i(s)\bigr)$ and $ p_0=\min_{i\in[r]}p_i(0).
$
By \eqref{eq:book-density-hyp}, $p_0\ge p>0$. At every stage at which the
construction is defined, put
\begin{equation}\label{eq:q-alpha}
 q_i(s)=p_i(s)-p_0+\delta,
 \qquad
 \alpha_i(s)=\frac{q_i(s)}{t}.
\end{equation}

We maintain the following structural invariants for every $i\in[r]$:
\begin{equation}\label{eq:book-invariants}
 \begin{aligned}
 &T_i(s)\subseteq X
   \quad\text{and}\quad T_i(s)\text{ is a color-$i$ clique},\\
 &X(s)\subseteq X\cap
   \bigcap_{h=1}^r\bigcap_{v\in T_h(s)}N_h(v),\qquad Y_i(s)\subseteq Y_i\cap\bigcap_{v\in T_i(s)}N_i(v).
 \end{aligned}
\end{equation}
The invariants hold at time $0$. Because the graph is simple, a vertex does
not belong to its own color-$i$ neighborhood. It follows from
\eqref{eq:book-invariants} that $X(s)$ is disjoint from every current spine
and that $T_i(s)\cap Y_i(s)=\emptyset$ for every $i$. Thus a vertex chosen
from $X(s)$ is always a genuinely new spine vertex.

Suppose that $|T_i(s)|<t$ for every $i\in[r]$ and that the current state is
defined with $p_i(s)>0$ and $q_i(s)>0$ for every $i$. Apply
\Cref{lem:key} to the current reservoir and page sets with parameters
$\alpha_i(s)$. We obtain a vertex $x\in X(s)$, a color $\ell\in[r]$, a
number $\lambda\ge-1$, a set $X'\subseteq X(s)$, and sets
$
 Y_i'\subseteq N_i(x)\cap Y_i(s)
~(i\in[r]).
$
We distinguish two cases.

If $\lambda\le\lambda_0$, then the edges from $x$ to
$X'\setminus\{x\}$ have one of the $r$ colors. Hence some $j\in[r]$
satisfies
\[
 |N_j(x)\cap X'|\ge\frac{|X'|-1}{r}.
\]
We extend the color-$j$ spine by setting
\[
 T_j(s+1)=T_j(s)\cup\{x\},
 \qquad
 X(s+1)=N_j(x)\cap X',
 \qquad
 Y_j(s+1)=Y_j',
\]
and set $T_i(s+1)=T_i(s)$ and $Y_i(s+1)=Y_i(s)$ for $i\ne j$.

If $\lambda>\lambda_0$, we perform a \emph{boost step} in color $\ell$ and
set
$
 X(s+1)=X'$
and $Y_\ell(s+1)=Y_\ell'$,
while leaving every spine and every other page set unchanged.

The structural invariants in \eqref{eq:book-invariants} are preserved. For a
boost step, the spines remain unchanged, $X(s+1)\subseteq X(s)$, and only
$Y_\ell(s)$ is replaced by a subset. For a spine-extension step in color
$j$, the induction hypothesis gives
\[
 x\in X(s)\subseteq\bigcap_{v\in T_j(s)}N_j(v),
\]
so $T_j(s)\cup\{x\}$ is a color-$j$ clique. Moreover,
\[
 X(s+1)\subseteq N_j(x)\cap X(s),
 \qquad
 Y_j(s+1)\subseteq N_j(x)\cap Y_j(s),
\]
and all other spines and page sets remain unchanged. Thus, in either type of
step,
\[
 X(s+1)\subseteq X(s),
 \qquad
 Y_i(s+1)\subseteq Y_i(s)
 \quad(i\in[r]).
\]
We next establish quantitative estimates showing, in particular, that none
of the sets or parameters needed above can vanish before a spine reaches
size $t$.

Formally, carry out the procedure only as long as the current reservoir is
nonempty, and stop at the first state in which some spine has size $t$. If a
proposed spine-extension step were to produce an empty new reservoir, stop
provisionally at that step. The estimates below apply to every legal finite
initial segment. The reservoir estimate will show that the provisional
stopping event is impossible.

\medskip
\noindent\emph{Density bookkeeping.}
For $i\in[r]$, let $B_i(s)$ be the set of earlier boost steps in color $i$,
and let $\lambda(q)$ denote the threshold used at step $q$. Also set
$ e_i(s)=|T_i(s)|$.
Restricting the first argument in a minimum relative degree cannot decrease
that degree. Thus, if the color-$j$ spine is extended at step $s$, then
\[
 p_j(s+1)\ge p_j(X',Y_j')\ge p_j(s)-\alpha_j(s),
\]
whereas, for $i\ne j$, the page set is unchanged and
$X(s+1)\subseteq X(s)$, so $p_i(s+1)\ge p_i(s)$. Similarly, if a boost step
is performed in color $\ell$, then
\[
 p_\ell(s+1)\ge p_\ell(s)+\lambda(s)\alpha_\ell(s),
 \qquad
 p_i(s+1)\ge p_i(s)\quad(i\ne\ell).
\]
Consequently, if the color-$i$ spine is extended at step $s$, then
\[
 q_i(s+1)
 \ge q_i(s)-\alpha_i(s)
 =\left(1-\frac1t\right)q_i(s).
\]
If a boost step is performed in color $i$, then
\[
 q_i(s+1)
 \ge q_i(s)+\lambda(s)\alpha_i(s)
 =\left(1+\frac{\lambda(s)}t\right)q_i(s).
\]
In every other case, $q_i(s+1)\ge q_i(s)$. Iterating these one-step
inequalities gives
\[
 q_i(s)
 \ge q_i(0)\left(1-\frac1t\right)^{e_i(s)}
 \prod_{q\in B_i(s)}\left(1+\frac{\lambda(q)}t\right).
\]
Since $q_i(0)\ge\delta$ and $e_i(s)\le t$ up to and including the first
time a spine reaches size $t$,
\begin{equation}\label{eq:q-product}
 q_i(s)
 \ge\delta\left(1-\frac1t\right)^t
 \prod_{q\in B_i(s)}\left(1+\frac{\lambda(q)}t\right).
\end{equation}
The assumptions imply $t\ge\lambda_0\ge2$, and hence
$(1-1/t)^t\ge1/4$. Therefore
\begin{equation}\label{eq:p-lower}
 q_i(s)\ge\frac\delta4,
 \qquad
 p_i(s)=p_0-\delta+q_i(s)
 \ge p_0-\frac{3\delta}{4}
 \ge p-\frac{3\delta}{4}>0.
\end{equation}
In particular, $q_i(s)>0$, so every $\alpha_i(s)$ in
\eqref{eq:q-alpha} is positive.

Since every relative density is at most $1$ and $\delta\le1/4$, we also have
$q_i(s)\le1+\delta\le5/4$. Combining this with \eqref{eq:q-product} yields
\[
 \prod_{q\in B_i(s)}\left(1+\frac{\lambda(q)}t\right)
 \le\frac5\delta.
\]
Because $\delta\le1/4$, we have $\log(5/\delta)\le3L$, and hence
\begin{equation}\label{eq:boost-log-sum}
 \sum_{q\in B_i(s)}
 \log\left(1+\frac{\lambda(q)}t\right)
 \le3L.
\end{equation}
At every boost step, $\lambda(q)>\lambda_0$. Since $t\ge\lambda_0$ and
$\log(1+x)\ge x/2$ for $0\le x\le1$, each summand in
\eqref{eq:boost-log-sum} is at least
\[
 \log\left(1+\frac{\lambda_0}{t}\right)
 \ge\frac{\lambda_0}{2t}.
\]
Consequently,
\begin{equation}\label{eq:boost-count}
 |B_i(s)|\le\frac{6L}{\lambda_0}t\le t.
\end{equation}

\medskip
\noindent\emph{Page-set retention.}
The page set $Y_i(s)$ is replaced only when the color-$i$ spine is extended
or when a boost step is performed in color $i$. On each such occasion,
\Cref{lem:key} gives
\[
 |Y_i(s+1)|=p_i(s)|Y_i(s)|.
\]
Set
$\xi=\frac{3\delta}{4p}.$
Then $0<\xi\le3/16$, and \eqref{eq:p-lower} gives
$
 p_i(s)\ge p-\frac{3\delta}{4}=p(1-\xi).
$
Let $b_i(s)=|B_i(s)|$ and put $u_i(s)=e_i(s)+b_i(s)$. Up to and including
the first state in which some spine reaches size $t$,
\[
 e_i(s)\le t,
 \qquad
 b_i(s)\le\frac{6L}{\lambda_0}t,
 \qquad
 u_i(s)\le2t.
\]
It follows that
$
 |Y_i(s)|
 \ge p^{t+6Lt/\lambda_0}(1-\xi)^{2t}|Y_i(0)|.
$
Since $\log(1-\xi)\ge-2\xi$ for $0\le\xi\le3/16$, we obtain
\begin{align}
 |Y_i(s)|
 &\ge p^t
 \exp\!\left(-\frac{6LL_p}{\lambda_0}t\right)
 \exp\!\left(-\frac{3\delta}{p}t\right)|Y_i(0)|=p^te^{-\Pi t}|Y_i(0)|.
 \label{eq:page-retention}
\end{align}
By \eqref{eq:book-page-hyp}, every page set occurring in the construction
therefore has size at least $m$.

\medskip
\noindent\emph{The cost of boost steps.}
Suppose that a boost step occurs at time $s$ in color $i$. Then
$
 p_i(s+1)\ge p_i(s)+\lambda(s)\alpha_i(s).
$
Since $p_i(s+1)\le1$ and, by \eqref{eq:p-lower},
$\alpha_i(s)=q_i(s)/t\ge\delta/(4t)$, we have
\begin{equation}\label{eq:lambda-upper}
 \lambda(s)\le\frac{4t}{\delta}.
\end{equation}
For every threshold $\lambda>\lambda_0$ arising from a boost step, we claim
that
\begin{equation}\label{eq:root-over-log}
 \frac{(\lambda+1)^{1/d}}
 {\log(1+\lambda/t)}
 \le\frac{4t}{\lambda_0^{(d-1)/d}}.
\end{equation}
If $\lambda\le t$, then $\lambda\ge\lambda_0\ge2$ and
$\log(1+\lambda/t)\ge\lambda/(2t)$. Hence
\[
 \frac{(\lambda+1)^{1/d}}{\log(1+\lambda/t)}
 \le2t\frac{(\lambda+1)^{1/d}}{\lambda}
 \le4t\lambda^{-(d-1)/d}
 \le\frac{4t}{\lambda_0^{(d-1)/d}}.
\]
If $\lambda>t$, then $\log(1+\lambda/t)\ge\log2$. Moreover,
\eqref{eq:lambda-upper}, $t\ge1$, and $\delta\le1$ give
$\lambda+1\le\frac{5t}{\delta}.
$
The hypothesis $t\ge\lambda_0\delta^{-1/(d-1)}$ is equivalent to
$t^{1/d}\delta^{-1/d}
 \le t\lambda_0^{-(d-1)/d}.
$
Since $d\ge3$ and $5^{1/d}/\log2<4$, inequality
\eqref{eq:root-over-log} follows in this case as well.

Multiplying \eqref{eq:root-over-log} by the corresponding logarithms, using
\eqref{eq:boost-log-sum}, and summing over all colors gives
\begin{equation}\label{eq:boost-root-sum}
 \sum_{q\in B(s)}(\lambda(q)+1)^{1/d}
 \le\frac{12rLt}{\lambda_0^{(d-1)/d}},
 \qquad
 B(s)=\bigcup_{i=1}^rB_i(s).
\end{equation}

\medskip
\noindent\emph{Reservoir retention.}
Set
\[
 \varepsilon_0=\frac\beta r
 \exp\!\left(-C(\lambda_0+1)^{1/d}\right).
\]
Then $0<\varepsilon_0\le1$. In a spine-extension step,
$\lambda\le\lambda_0$, so \Cref{lem:key} and the choice of the majority
color give
\begin{align*}
 |X(s+1)|
 \ge\frac{|X'|-1}{r}\ge\frac\beta r
       e^{-C(\lambda_0+1)^{1/d}}|X(s)|-\frac1r\ge\varepsilon_0|X(s)|-1.
\end{align*}
In a boost step, \Cref{lem:key} gives
$
 |X(s+1)|
 \ge\beta e^{-C(\lambda(s)+1)^{1/d}}|X(s)|.
$
Since $\varepsilon_0\le\beta$, this implies
$
 |X(s+1)|
\ge\varepsilon_0e^{-C(\lambda(s)+1)^{1/d}}|X(s)|.
$
For each step $q$, define
\[
 a_q=
 \begin{cases}
  \varepsilon_0,&\text{if a spine is extended at step $q$},\\
  \varepsilon_0e^{-C(\lambda(q)+1)^{1/d}},
  &\text{if a boost step is performed at step $q$},
 \end{cases}
\]
and let $\epsilon_q=1$ in the first case and $\epsilon_q=0$ in the second.
With $R_q=|X(q)|$, both estimates take the form
$
 R_{q+1}\ge a_qR_q-\epsilon_q,$ where $
 0<a_q\le1.
$
Before or at the first state in which some spine reaches size $t$, fewer than
$rt$ spine extensions and at most $rt$ boost steps have occurred. Indeed, at
such a state the number of spine extensions is at most
$
 t+(r-1)(t-1)=rt-r+1<rt,
$
and \eqref{eq:boost-count} bounds the number of boost steps by $rt$. Thus
there are at most $2rt$ factors $a_q$. Furthermore,
$\log(1/\varepsilon_0)
=\rho+C(\lambda_0+1)^{1/d}
\le\rho+2C\lambda_0^{1/d}.
$
Together with \eqref{eq:boost-root-sum}, this gives, for every constructed
initial segment and also for a putative final spine-extension step that would
empty the reservoir,
\begin{align*}
 \prod_{q=0}^{s-1}a_q
 &=\varepsilon_0^s
 \exp\!\left(-C\sum_{q\in B(s)}(\lambda(q)+1)^{1/d}\right)\ge\exp\!\left(
-2rt\bigl(\rho+2C\lambda_0^{1/d}\bigr)
 -\frac{12CrLt}{\lambda_0^{(d-1)/d}}
 \right)\ge e^{-rt\Xi}.
\end{align*}
Iterating the recurrence for $R_q$ yields
\[
 R_s\ge
 \left(\prod_{q=0}^{s-1}a_q\right)R_0
 -\sum_{j=0}^{s-1}\epsilon_j
  \prod_{q=j+1}^{s-1}a_q.
\]
Every product in the sum is at most $1$, and fewer than $rt$ of the
$\epsilon_j$ are nonzero. Therefore, by
\eqref{eq:book-reservoir-hyp},
$
 |X(s)|=R_s
 \ge e^{-rt\Xi}|X(0)|-rt
 \ge rt>0.
$

We may now close the induction implicit in the construction. The last
reservoir estimate rules out a proposed spine extension with an empty new
reservoir. At every stage before a spine reaches size $t$,
\eqref{eq:p-lower} gives $p_i(s)>0$ and $q_i(s)>0$,
\eqref{eq:page-retention} gives $|Y_i(s)|\ge m\ge1$, and the reservoir
estimate gives $|X(s)|>0$. Therefore \Cref{lem:key} is always applicable,
and the construction cannot stop prematurely.

By \eqref{eq:boost-count}, at most $rt$ boost steps can occur in total. If
every spine had size at most $t-1$, fewer than $rt$ spine extensions could
occur. Since every step is of one of these two types, the process cannot
continue indefinitely without producing a spine of size $t$. Hence, for some
color $i$, the construction reaches a state $s$ with $|T_i(s)|=t$.

Set $T=T_i(s)$. By \eqref{eq:book-invariants}, $T\subseteq X$ is a
color-$i$ clique, and
\[
 Y_i(s)\subseteq Y_i\cap\bigcap_{v\in T}N_i(v).
\]
The page-retention estimate \eqref{eq:page-retention} gives
$|Y_i(s)|\ge m$. Choose a set $P\subseteq Y_i(s)$ with $|P|=m$. Since the
graph is simple, the preceding inclusion implies $Y_i(s)\cap T=\emptyset$.
Consequently,
$
 T\subseteq X$ and
 $P\subseteq Y_i\setminus T,$
and every edge between $T$ and $P$ has color $i$. Thus $(T,P)$ has all the
properties asserted in the theorem.
\end{proof}

\section{Proof of the Ramsey bounds}\label{sec:transfer}

\subsection{Retained spines and an entropy estimate}\label{sec:spines}

\begin{lemma}\label{lem:spine-compatibility}
Let $r,k\ge2$ be integers, and let $S_1,\ldots,S_r,W$ be as in
\Cref{lem:regularization}. Write
$
 s_j=|S_j|
~(j\in[r]).
$
Fix $i\in[r]$, and let $(T,P)$ be a color-$i$ book contained in $W$; thus
$T\cup P\subseteq W$.

If $s_i+|T|\ge k$, then $S_i\cup T$ contains a color-$i$ copy of $K_k$. If
$s_i+|T|<k$ and $P$ contains a color-$i$ clique of size
$k-s_i-|T|$, then that clique together with $S_i\cup T$ forms a color-$i$
copy of $K_k$.

Moreover, let $j\ne i$. If $s_j\ge k$, then $S_j$ contains a color-$j$ copy
of $K_k$. If $s_j<k$ and $P$ contains a color-$j$ clique of size $k-s_j$,
then that clique together with $S_j$ forms a color-$j$ copy of $K_k$.
\end{lemma}

\begin{proof}
By \Cref{lem:regularization}, for every $j\in[r]$, the set $S_j$ is a
color-$j$ clique, $S_j\cap W=\emptyset$, and every edge between $S_j$ and
$W$ has color $j$. Since $(T,P)$ is a color-$i$ book, the sets $T$ and $P$
are disjoint, $T$ is a color-$i$ clique, and every edge between $T$ and $P$
has color $i$.

Because $T\subseteq W$, every edge between $S_i$ and $T$ has color $i$.
It follows that $S_i\cup T$ is a color-$i$ clique of size $s_i+|T|$. Hence,
if $s_i+|T|\ge k$, this clique contains a color-$i$ copy of $K_k$.
Now suppose that $s_i+|T|<k$, and let $Q\subseteq P$ be a color-$i$ clique
of size
$
 |Q|=k-s_i-|T|.
$
The sets $S_i,T,Q$ are pairwise disjoint. The edges inside each of these sets
have color $i$; every edge between $S_i$ and $T\cup Q$ has color $i$ because
$T\cup Q\subseteq W$; and every edge between $T$ and $Q$ has color $i$ by
the book property. Thus $S_i\cup T\cup Q$ is a color-$i$ clique of size
$s_i+|T|+|Q|=k.
$

Finally, let $j\ne i$. If $s_j\ge k$, then the color-$j$ clique $S_j$
contains a color-$j$ copy of $K_k$. Suppose instead that $s_j<k$, and let
$Q\subseteq P$ be a color-$j$ clique of size $k-s_j$. Since $Q\subseteq W$,
the sets $S_j$ and $Q$ are disjoint, and every edge between them has color
$j$. Therefore $S_j\cup Q$ is a color-$j$ clique of size $k$.
\end{proof}

\begin{lemma}\label{lem:entropy-loss}
Let $r,k\ge2$ and $t\ge1$ be integers. Let $s_1,\ldots,s_r$ be nonnegative
integers satisfying
\[
 s_j<k\quad(j\in[r]),
 \qquad
 s=\sum_{j=1}^rs_j<\frac{rt}{4}.
\]
Fix $i\in[r]$ such that $s_i+t<k$, and put
\[
 b_j=
 \begin{cases}
  k-s_i-t,&j=i,\\
  k-s_j,&j\ne i,
 \end{cases}
 \qquad
 B=\sum_{j=1}^rb_j=rk-s-t.
\]
Then
\begin{equation}\label{eq:entropy-loss}
 R(b_1,\ldots,b_r)
 \le r^{rk-s-t}\exp\!\left(-\frac{t^2}{64k}\right).
\end{equation}
\end{lemma}

\begin{proof}
The assumptions $s_j<k$ for all $j$ and $s_i+t<k$ imply that
$b_1,\ldots,b_r$ are positive integers. By
\eqref{eq:off-diagonal-basic},
\begin{equation}\label{eq:entropy-multinomial-bound}
 R(b_1,\ldots,b_r)
 \le\binom{B}{b_1,\ldots,b_r}.
\end{equation}

Let $(M_1,\ldots,M_r)$ have the multinomial distribution with $B$ trials and
equal cell probabilities $1/r$. Then
\begin{equation}\label{eq:multinomial-coordinate}
 r^{-B}\binom{B}{b_1,\ldots,b_r}
 =\Pp(M_1=b_1,\ldots,M_r=b_r)
 \le\Pp(M_i=b_i).
\end{equation}
Set
$
 q=\frac1r$ and
$ x=\frac{b_i}{B}$.
The random variable $M_i$ has distribution $\operatorname{Bin}(B,q)$, and
$xB=b_i$ is an integer. Hence
\[
 \Pp(M_i=b_i)
 =\binom{B}{xB}q^{xB}(1-q)^{(1-x)B}.
\]
Using
\[
 \binom{B}{xB}\le\exp\bigl(BH(x)\bigr),
 \qquad
 H(x)=-x\log x-(1-x)\log(1-x),
\]
we obtain
\begin{equation}\label{eq:binomial-entropy}
 \Pp(M_i=b_i)\le\exp\bigl(-B\KL(x\|q)\bigr),
\end{equation}
where
\[
 \KL(x\|q)
 =x\log\frac{x}{q}+(1-x)\log\frac{1-x}{1-q}.
\]

We next quantify how far the distinguished coordinate lies below its
multinomial mean. Define
\[
 \Delta=\frac Br-b_i.
\]
Since $B=rk-s-t$ and $b_i=k-s_i-t$,
\[
 \Delta
 =s_i+\left(1-\frac1r\right)t-\frac sr
 >\left(1-\frac1r\right)t-\frac t4
 \ge\frac t4.
\]
Thus $b_i<B/r$, so $0<x<q$, and
\begin{equation}\label{eq:entropy-displacement}
 q-x=\frac{\Delta}{B}\ge\frac{t}{4B}.
\end{equation}

For fixed $q$, let
$
 \phi(u)=\KL(u\|q)
 ~(0<u<1).
$
Then
\[
 \phi(q)=\phi'(q)=0,
 \qquad
 \phi''(u)=\frac1{u(1-u)}.
\]
For every $u\in[x,q]$, we have $u\le q=1/r$, and therefore
\[
 \phi''(u)\ge\frac1u\ge r\ge\frac r2.
\]
Taylor's theorem gives
\[
 \KL(x\|q)=\phi(x)\ge\frac r4(x-q)^2.
\]
Combining this bound with \eqref{eq:entropy-displacement}, we find
\[
 B\KL(x\|q)
 \ge\frac r{4B}\Delta^2
 \ge\frac{rt^2}{64B}
 \ge\frac{t^2}{64k},
\]
where the last inequality uses $B=rk-s-t\le rk$. Finally,
\eqref{eq:entropy-multinomial-bound},
\eqref{eq:multinomial-coordinate}, and \eqref{eq:binomial-entropy} give
\[
 R(b_1,\ldots,b_r)
 \le r^B\exp\!\left(-\frac{t^2}{64k}\right)
 =r^{rk-s-t}\exp\!\left(-\frac{t^2}{64k}\right).
\]
This proves \eqref{eq:entropy-loss}.
\end{proof}

\subsection{Choice of parameters and completion of the proof}
\label{sec:transfer-theorem}

For $b\in(0,1)$ and integers $r\ge2$ and $d\ge3$, define
\begin{equation}\label{eq:vartheta}
 \vartheta_{r,d}(b)=
 \frac{b}{
 r^{d/(d-1)}d^{2d/(d-1)}
 \bigl(\log(2rd)\bigr)^{2/(d-1)}}.
\end{equation}
Once $b$ has been fixed, we abbreviate
$\vartheta_{r,d}=\vartheta_{r,d}(b)$.

\begin{theorem}\label{thm:transfer}
There exist absolute constants $b\in(0,1)$ and $c,K>0$ such that the
following holds. Let $r,k\ge2$ and $d\ge3$ be integers satisfying
\[
 d\le\max\{3,\log(2r)\}.
\]
If
\begin{equation}\label{eq:transfer-threshold}
 k\ge
 K r^{(2d^2-1)/(d-1)^2}
 d^{2d(2d-1)/(d-1)^2}
 \bigl(\log(2rd)\bigr)^{2d^2/(d-1)^2},
\end{equation}
then, with $\vartheta_{r,d}=\vartheta_{r,d}(b)$,
\begin{equation}\label{eq:transfer-saving}
 R_r(k)\le
 \exp\!\left(-c\vartheta_{r,d}^2k\right)r^{rk}.
\end{equation}
\end{theorem}

We now choose the absolute constants used in the proof. First take $A\ge12$
sufficiently large. Then choose $0<\alpha,\zeta\le1/8$, and finally choose
$0<\gamma\le\min\{1/8,\alpha/64\}$, so that
\begin{equation}\label{eq:constant-choice}
 4\alpha+3\zeta+\frac6A+2\gamma<\frac1{256}.
\end{equation}
For admissible $r,d,k$, write
\[
 \ell=\log(2rd),
 \qquad
 \vartheta=\vartheta_{r,d}(b),
 \qquad
 \tau=\frac{d-1}{d},
\]
and define
\begin{gather}
 t=\lfloor\vartheta k\rfloor,
 \qquad
 \eta=\frac{\alpha\vartheta}{r},
 \qquad
 p=\frac1r-2\eta=\frac{1-2\alpha\vartheta}{r},
 \notag\\
 \delta=\zeta p\vartheta,
 \qquad
 L=\log(1/\delta),
 \qquad
 L_p=\log(1/p),
 \qquad
 \lambda_0=\frac{AL^2}{\vartheta},
 \qquad
 \Gamma=\gamma\vartheta^2.
 \label{eq:transfer-lambda-gamma}
\end{gather}
For the correlation constants in \eqref{eq:correlation-constants}, put
\begin{equation}\label{eq:transfer-rho-Pi-Xi}
 \rho=\log(r/\beta_r),
 \qquad
 \Pi=\frac{3\delta}{p}+\frac{6LL_p}{\lambda_0},
 \qquad
 \Xi=2\rho+4C_{r,d}\lambda_0^{1/d}
 +\frac{12C_{r,d}L}{\lambda_0^\tau}.
\end{equation}

The next lemma fixes $b$ and verifies, uniformly in all the remaining
parameters, the hypotheses needed in the book argument.

\begin{lemma}\label{lem:uniform-parameters}
There are absolute choices of $b\in(0,1)$ and $K_0\ge1$ such that the
following holds. Let $r\ge2$, let $d$ be an integer satisfying
$
 3\le d\le\max\{3,\log(2r)\},
$
and let $k$ be an integer satisfying
\[
 k\ge
 K_0 r^{(2d^2-1)/(d-1)^2}
 d^{2d(2d-1)/(d-1)^2}
 \bigl(\log(2rd)\bigr)^{2d^2/(d-1)^2}.
\]
Then all of the following statements hold.
\begin{enumerate}[label=\textup{(\roman*)}]
\item The basic parameters satisfy
\begin{equation}\label{eq:uniform-basic}
 0<\vartheta\le\frac18,
 \qquad
 \frac1{2r}\le p\le\frac1r,
 \qquad
 0<\delta\le\min\{p/4,1/4\},
 \qquad
 \eta\le1,
\end{equation}
and
\begin{equation}\label{eq:uniform-logs}
 c_{\log}\ell\le L_p\le L\le C_{\log}\ell,
 \qquad
 \lambda_0\ge\max\{2,6L\},
\end{equation}
for absolute constants $c_{\log},C_{\log}>0$.

\item The integer $t$ satisfies
\begin{equation}\label{eq:uniform-book}
 \frac{\vartheta k}{2}\le t\le\vartheta k,
 \qquad
 t\le\frac k8,
 \qquad
 t\ge\lambda_0\delta^{-1/(d-1)},
\end{equation}
and
\begin{equation}\label{eq:uniform-margins}
 \left(\frac{\alpha}{16}-\gamma\right)\vartheta^2k\ge\log2,
 \qquad
 \frac{\vartheta t}{256}\ge\log2.
\end{equation}

\item We have
\begin{equation}\label{eq:uniform-size}
 e^{-\Gamma k}r^{rk}\ge2,
 \qquad
 \Gamma k+\log2\le\frac1{32}rk\log r,
\end{equation}
and
\begin{equation}\label{eq:uniform-auxiliary}
 \log(1/\eta)\le\frac34rk\log r,
 \qquad
 \log(2rt)\le\frac18rk\log r.
\end{equation}

\item The total reservoir loss satisfies
\begin{equation}\label{eq:uniform-reservoir-exponent}
 rt\Xi\le\frac18rk\log r.
\end{equation}
\end{enumerate}
\end{lemma}

\begin{proof}
Set $\ell_r=\log(2r)$. The restriction on $d$ implies that
\begin{equation}\label{eq:ell-log-comparison}
 \ell_r\le\ell=\log(2rd)\le C_{\ell}\ell_r
\end{equation}
for an absolute constant $C_{\ell}>0$. Indeed, if $\ell_r<3$, then $d=3$,
and the assertion follows after increasing $C_{\ell}$. If $\ell_r\ge3$,
then $d\le\ell_r$, and hence
$
 \ell=\ell_r+\log d
 \le\ell_r+\log\ell_r
 \le2\ell_r.
$
Since $r\ge2$, it also follows that $\ell\le C\log r$ for an absolute
constant $C>0$.

We first require $b\in(0,1)$ to be small enough that
\[
 \vartheta\le\frac18,
 \qquad
 2\alpha\vartheta\le\frac12,
 \qquad
 \zeta\vartheta\le\frac14
\]
for every admissible pair $(r,d)$. Then
\[
 \frac1{2r}\le
 p=\frac{1-2\alpha\vartheta}{r}
 \le\frac1r,
\]
and
$
 0<\delta=\zeta p\vartheta
 \le\min\{p/4,1/4\}.
$
Also, $\eta=\alpha\vartheta/r\le1$. Thus
\eqref{eq:uniform-basic} holds.

Put $B_b=1+\log(1/b)$. Since
$1/2\le1-2\alpha\vartheta\le1$,
\[
 \log r\le L_p
 =\log r-\log(1-2\alpha\vartheta)
 \le\log r+\log2.
\]
Moreover,
$L=\log(1/\zeta)+L_p+\log(1/\vartheta),
$
where
\[
 \log(1/\vartheta)
 =\log(1/b)
 +\frac{d}{d-1}\log r
 +\frac{2d}{d-1}\log d
 +\frac{2}{d-1}\log\ell.
\]
Because $d\ge3$, all the coefficients in the last display are bounded by
absolute constants. It follows from \eqref{eq:ell-log-comparison} that
\begin{equation}\label{eq:uniform-log-prechoice}
 c\ell\le L_p\le L\le CB_b\ell
\end{equation}
for absolute constants $c,C>0$.

Since $\delta\le1/4$, we have $L\ge\log4$. Therefore
\[
 \frac{\lambda_0}{L}
 =\frac{AL}{\vartheta}
 \ge8A\log4>6,
\]
and $\lambda_0\ge2$. Once $b$ has been fixed below,
\eqref{eq:uniform-log-prechoice} gives \eqref{eq:uniform-logs} for suitable
absolute constants $c_{\log},C_{\log}>0$.

We next determine the scale required by the book condition. Since
$p\ge1/(2r)$ and $d\ge3$,
\[
 \lambda_0\delta^{-1/(d-1)}
 =AL^2\vartheta^{-1}(\zeta p\vartheta)^{-1/(d-1)}
 \le CA L^2r^{1/(d-1)}\vartheta^{-d/(d-1)},
\]
where $C>0$ is absolute because $\zeta$ has already been fixed. Hence it is
enough to require
\begin{equation}\label{eq:uniform-threshold-reduction}
 k\ge CA L^2r^{1/(d-1)}
 \vartheta^{-(2d-1)/(d-1)}.
\end{equation}
Indeed, after increasing $C$ if necessary, this inequality gives
\[
 \frac{\vartheta k}{2}
 \ge\lambda_0\delta^{-1/(d-1)}.
\]
The quantity on the right is at least $\lambda_0\ge2$, so
$\vartheta k\ge4$. Consequently,
\[
 t=\lfloor\vartheta k\rfloor
 \ge\frac{\vartheta k}{2}
 \ge\lambda_0\delta^{-1/(d-1)}.
\]

Using \eqref{eq:uniform-log-prechoice} and substituting
\eqref{eq:vartheta} into \eqref{eq:uniform-threshold-reduction}, its
right-hand side is at most
\[
 C A B_b^2 b^{-(2d-1)/(d-1)}
 r^{(2d^2-1)/(d-1)^2}
 d^{2d(2d-1)/(d-1)^2}
 \ell^{2d^2/(d-1)^2}.
\]
Here
\[
 \frac1{d-1}+\frac{d(2d-1)}{(d-1)^2}
 =\frac{2d^2-1}{(d-1)^2}
\]
and
\[
 2+\frac{2(2d-1)}{(d-1)^2}
 =\frac{2d^2}{(d-1)^2}.
\]
Since
\[
 2\le\frac{2d-1}{d-1}\le\frac52
 \qquad(d\ge3),
\]
the power of $1/b$ is uniformly bounded once $b$ is fixed. Thus the
threshold in the statement of the lemma, with a sufficiently large absolute
constant after $b$ has been fixed, implies
\eqref{eq:uniform-threshold-reduction}.

We now choose $b$ by estimating the reservoir loss. By
\eqref{eq:correlation-constant-bounds}, there are absolute constants
$C_{\mathrm{corr}},C_{\rho}>0$ such that
\begin{equation}\label{eq:uniform-C-rho}
 C_{r,d}\le C_{\mathrm{corr}}rd^2\ell,
 \qquad
 \rho=\log r+\log(1/\beta_r)\le C_{\rho}r.
\end{equation}
The definition of $\vartheta$ gives the exact identity
\begin{equation}\label{eq:uniform-balancing}
 \vartheta^\tau rd^2\ell^{2/d}=b^\tau.
\end{equation}
Using $t\le\vartheta k$, \eqref{eq:transfer-rho-Pi-Xi},
\eqref{eq:transfer-lambda-gamma}, \eqref{eq:uniform-C-rho},
\eqref{eq:uniform-log-prechoice}, and \eqref{eq:uniform-balancing}, we first
obtain
$
 2rt\rho\le Cr^2\vartheta k.
$
For the second contribution to $rt\Xi$,
\begin{align*}
 4rtC_{r,d}\lambda_0^{1/d}
 \le C A^{1/d}r^2kd^2\ell\,
 \vartheta^\tau L^{2/d}=C A^{1/d}b^\tau rk\,
 \ell^{1-2/d}L^{2/d}\le C A^{1/d}b^\tau B_b^{2/d}rk\ell.
\end{align*}
For the third contribution,
\begin{align*}
 \frac{12rtC_{r,d}L}{\lambda_0^\tau}
 \le C A^{-\tau}r^2kd^2\ell\,
 \vartheta^{1+\tau}L^{1-2\tau}=C A^{-\tau}b^\tau r\vartheta k
 \left(\frac{\ell}{L}\right)^{(d-2)/d}\le C A^{-\tau}b^\tau r\vartheta k,
\end{align*}
where the last inequality uses $L\ge c\ell$.

Since
\[
 r\vartheta\le b,
 \qquad
 \tau\ge\frac23,
 \qquad
 A^{1/d}\le A^{1/3},
 \qquad
 B_b^{2/d}\le B_b^{2/3},
\]
and $\ell\le C\log r$, division by $rk\log r$ gives
\[
 \frac{rt\Xi}{rk\log r}
 \le C\left(
 b+A^{1/3}b^{2/3}B_b^{2/3}
 +A^{-2/3}b^{2/3}
 \right).
\]
The expression on the right tends to $0$ as $b\downarrow0$. We may therefore
fix an absolute $b\in(0,1)$ so small that this expression is at most $1/8$
and all the earlier smallness conditions hold. This proves
\eqref{eq:uniform-reservoir-exponent}. The choice of $b$ is independent of
$r,d,k$. With this value of $b$ fixed,
\eqref{eq:uniform-log-prechoice} also gives \eqref{eq:uniform-logs} for
suitable absolute constants $c_{\log},C_{\log}>0$.

It remains to choose $K_0$. Define
\[
 \mathcal T_{r,d}=
 r^{(2d^2-1)/(d-1)^2}
 d^{2d(2d-1)/(d-1)^2}
 \ell^{2d^2/(d-1)^2}.
\]
A direct calculation gives
\[
 \vartheta^2\mathcal T_{r,d}
 =b^2
 r^{(2d-1)/(d-1)^2}
 d^{2d/(d-1)^2}
 \ell^{2(d^2-2d+2)/(d-1)^2}.
\]
All three bases on the right are greater than $1$, and all three exponents
are positive. Therefore, whenever $k\ge K_0\mathcal T_{r,d}$,
$
 \vartheta^2k\ge K_0b^2.
$

First choose $K_0$ large enough that
\eqref{eq:uniform-threshold-reduction} holds. Increase it further so that
\[
 \left(\frac{\alpha}{16}-\gamma\right)\vartheta^2k\ge\log2;
\]
this is possible because
$\alpha/16-\gamma\ge3\alpha/64>0$. The floor estimate already proved gives
\[
 \frac{\vartheta t}{256}
 \ge\frac{\vartheta^2k}{512},
\]
so another absolute increase of $K_0$ gives the second inequality in
\eqref{eq:uniform-margins}. Since $\vartheta\le1/8$, we also have
$t\le\vartheta k\le k/8$. We have now proved
\eqref{eq:uniform-book} and \eqref{eq:uniform-margins}.

Set
$
 H=rk\log r.
$
Since $\gamma\le1/8$, $\vartheta\le1/8$, and $r\log r\ge1$, we have
\[
 \Gamma k=\gamma\vartheta^2k\le\frac{H}{64}.
\]
By increasing $K_0$ so that $k$ is sufficiently large, we also have
$\log2\le H/64$. Hence
\[
 \Gamma k+\log2\le\frac{H}{32}.
\]
This is the second inequality in \eqref{eq:uniform-size}, and it also gives
$
 \log\bigl(e^{-\Gamma k}r^{rk}\bigr)
 =H-\Gamma k\ge\log2,
$
which proves the first one.

Finally, after $b$ has been fixed,
$ \log(1/\eta)
=\log(r/\alpha)+\log(1/\vartheta)
 \le C\ell\le C'\log r
$
for absolute constants $C,C'>0$. A further increase of $K_0$ gives the first
inequality in \eqref{eq:uniform-auxiliary}. Moreover, $t\le k$, and for all
sufficiently large $k$,
\[
 \log(2rt)
 \le\log(2rk)
 =\log(2r)+\log k
 \le\frac18rk\log r.
\]
For example, this follows uniformly once $k\ge128$ from
$\log(2r)\le r\log r$ and $\log k\le k/16$. Since
$\vartheta^2k\ge K_0b^2$, this lower bound on $k$ is ensured by one final
absolute increase of $K_0$. This proves \eqref{eq:uniform-auxiliary} and
completes the proof.
\end{proof}

\begin{proof}[Proof of \Cref{thm:transfer}]
Let $b$ and $K_0$ be the absolute constants supplied by
\Cref{lem:uniform-parameters}, and take $K=K_0$. Set
\begin{equation}\label{eq:N-floor}
 N=\left\lfloor e^{-\Gamma k}r^{rk}\right\rfloor.
\end{equation}
By \eqref{eq:uniform-size}, the quantity inside the floor is at least $2$.
Thus
\begin{equation}\label{eq:N-lower}
 N\ge e^{-\Gamma k}r^{rk}-1
 \ge\frac12e^{-\Gamma k}r^{rk}.
\end{equation}

Consider an arbitrary $r$-coloring of $K_N$, and apply
\Cref{lem:regularization} with parameter $\eta$. Let
\[
 s_i=|S_i|,
 \qquad
 s=\sum_{i=1}^rs_i.
\]
If some $s_i\ge k$, then the color-$i$ clique $S_i$ already contains a
monochromatic copy of $K_k$. We may therefore assume that
$
 s_i<k
~(i\in[r]).
$

Suppose first that $s\ge rt/4$. Since $\eta\le1$, we have
$\log(1+\eta)\ge\eta/2$, and hence
\[
 s\log(1+\eta)
 \ge\frac{rt}{4}\cdot\frac\eta2
 =\frac{\alpha\vartheta t}{8}
 \ge\frac{\alpha}{16}\vartheta^2k
 \ge\Gamma k+\log2.
\]
The last two inequalities follow from \eqref{eq:uniform-book} and
\eqref{eq:uniform-margins}. Combining
\eqref{eq:regularization-size} and \eqref{eq:N-lower}, we obtain
\[
 |W|
 \ge\frac12e^{-\Gamma k}(1+\eta)^sr^{rk-s}
 \ge r^{rk-s}.
\]
For $j\in[r]$, the integer $k-s_j$ is positive, and
\eqref{eq:off-diagonal-basic} gives
$
 R(k-s_1,\ldots,k-s_r)\le r^{rk-s}.
$
Consequently, $W$ contains a color-$j$ clique of size $k-s_j$ for some
$j\in[r]$. Since $S_j\cap W=\emptyset$ and every edge between $S_j$ and
$W$ has color $j$, this clique together with $S_j$ forms a color-$j$ copy of
$K_k$.

We may therefore assume that
\begin{equation}\label{eq:small-spine}
 s<\frac{rt}{4}.
\end{equation}
Let
$
 I=\{i\in[r]:s_i+t<k\}.
$
This set is nonempty. Indeed, some $i\in[r]$ satisfies
$s_i\le s/r<t/4$, and hence
\[
 s_i+t<\frac{5t}{4}\le\frac{5k}{32}<k,
\]
where we used $t\le k/8$. For $i\in I$, define
\[
 m_i=R(k-s_1,\ldots,k-s_i-t,\ldots,k-s_r),
 \qquad
 m=\max_{i\in I}m_i.
\]
All Ramsey parameters in the definition of $m_i$ are positive integers.
Applying \Cref{lem:entropy-loss} separately to every $i\in I$ and taking
the maximum gives
\begin{equation}\label{eq:m-bound}
 m\le r^{rk-s-t}
 \exp\!\left(-\frac{t^2}{64k}\right).
\end{equation}

We next verify the hypotheses of the book lemma inside $W$. From
\eqref{eq:regularization-size}, \eqref{eq:N-lower},
\eqref{eq:small-spine}, and \eqref{eq:uniform-size},
\begin{align}
 \log|W|
 &\ge rk\log r-s\log r-\Gamma k-\log2\notag\\
 &\ge rk\log r-\frac{r\vartheta k}{4}\log r
       -\frac1{32}rk\log r\notag\\
 &\ge\frac{15}{16}rk\log r
 \ge\frac34rk\log r.
 \label{eq:W-crude}
\end{align}
Together with \eqref{eq:uniform-auxiliary}, this implies
$\log|W|\ge\log(1/\eta)$, or equivalently $\eta|W|\ge1$. Therefore, for
every $w\in W$ and $i\in[r]$, \eqref{eq:regularization-degree} gives
\[
 |N_i(w)\cap W|
 \ge\left(\frac1r-\eta\right)|W|-1
 \ge\left(\frac1r-2\eta\right)|W|
 =p|W|.
\]

By \Cref{thm:correlation}, the property
$\mathcal G_r^{(d)}(\beta_r,C_{r,d})$ holds. We apply \Cref{thm:book} with
$
 X=Y_1=\cdots=Y_r=W,
$
with correlation parameters $(\beta_r,C_{r,d})$, and with $\Pi,\Xi$ from
\eqref{eq:transfer-rho-Pi-Xi}. The statement of \Cref{thm:book} does not
require $X,Y_1,\ldots,Y_r$ to be pairwise disjoint. The preceding degree
estimate verifies \eqref{eq:book-density-hyp}, and the remaining parameter
assumptions of the book theorem follow from
\eqref{eq:uniform-basic}--\eqref{eq:uniform-book}.

It remains to verify the page and reservoir size assumptions. First,
\[
 \Pi
 =3\zeta\vartheta+\frac{6\vartheta L_p}{AL}
 \le\left(3\zeta+\frac6A\right)\vartheta.
\]
Since $2\alpha\vartheta\le1/2$, the inequality
$\log(1-x)\ge-2x$ for $0\le x\le1/2$ gives
\[
 \log(pr)=\log(1-2\alpha\vartheta)
 \ge-4\alpha\vartheta.
\]
Moreover, \eqref{eq:uniform-book} gives
\[
 \frac{t^2}{64k}\ge\frac{\vartheta t}{128},
 \qquad
 \Gamma k\le2\gamma\vartheta t.
\]
Using \eqref{eq:regularization-size}, \eqref{eq:N-lower}, and
\eqref{eq:m-bound}, and discarding the nonnegative term
$s\log(1+\eta)$, we obtain
\begin{align*}
 \log\frac{|W|}{p^{-t}e^{\Pi t}m}
 &\ge t\log(pr)-\Pi t+\frac{t^2}{64k}
       -\Gamma k-\log2\\
 &\ge\left(
 \frac1{128}-4\alpha-3\zeta-\frac6A-2\gamma
 \right)\vartheta t-\log2\\
 &\ge\frac{\vartheta t}{256}-\log2
 \ge0,
\end{align*}
where the last line uses \eqref{eq:constant-choice} and
\eqref{eq:uniform-margins}. Hence
$
 |W|\ge p^{-t}e^{\Pi t}m.
$

Finally, \eqref{eq:W-crude}, \eqref{eq:uniform-auxiliary}, and
\eqref{eq:uniform-reservoir-exponent} give
\[
 \log(2rt)+rt\Xi
 \le\frac14rk\log r
 \le\log|W|.
\]
Therefore
$
 |W|\ge2rt\,e^{rt\Xi}.
$
All the hypotheses of \Cref{thm:book} are now satisfied.

By \Cref{thm:book}, there exist an index $i\in[r]$, a color-$i$ clique
$T\subseteq W$ with $|T|=t$, and a set
$
 P\subseteq W\setminus T$ with
 $|P|=m$,
such that every edge between $T$ and $P$ has color $i$.

If $i\notin I$, then $s_i+t\ge k$. Since $T\subseteq W$,
\Cref{lem:spine-compatibility} shows that $S_i\cup T$ contains a color-$i$
copy of $K_k$. Suppose instead that $i\in I$. Since $|P|=m\ge m_i$, the
definition of $m_i$ implies that the coloring induced by $P$ contains either
a color-$i$ clique of size $k-s_i-t$, or a color-$j$ clique of size
$k-s_j$ for some $j\ne i$. In either case,
\Cref{lem:spine-compatibility} gives a monochromatic copy of $K_k$.

We have proved that every $r$-coloring of $K_N$ contains a monochromatic
$K_k$. Hence, by the definition of the Ramsey number and
\eqref{eq:N-floor},
\[
 R_r(k)\le N
 \le e^{-\Gamma k}r^{rk}
 =\exp(-\gamma\vartheta^2k)r^{rk}.
\]
This proves \eqref{eq:transfer-saving} with $c=\gamma$.
\end{proof}

\begin{proof}[Proof of \Cref{thm:parameterised-main}]
Let $b,c_{\mathrm{tr}},K_{\mathrm{tr}}>0$ be the absolute constants supplied
by \Cref{thm:transfer}. For every admissible pair $(r,d)$,
\[
 \vartheta_{r,d}^2
 =\frac{b^2}{
 r^{2d/(d-1)}d^{4d/(d-1)}
 \bigl(\log(2rd)\bigr)^{4/(d-1)}}.
\]
The hypothesis in \eqref{eq:parameter-threshold}, with
$K=K_{\mathrm{tr}}$, is exactly the threshold required in
\Cref{thm:transfer}. That theorem therefore gives
\[
 R_r(k)\le
 \exp\!\left(
 -c_{\mathrm{tr}}b^2\frac{k}{
 r^{2d/(d-1)}d^{4d/(d-1)}
 \bigl(\log(2rd)\bigr)^{4/(d-1)}}
 \right)r^{rk}.
\]
Thus \eqref{eq:parameter-saving} holds with the absolute constants
$c=c_{\mathrm{tr}}b^2$ and $K=K_{\mathrm{tr}}$.
\end{proof}

\begin{proof}[Proof of \Cref{thm:main}]
Set
\[
 \ell_r=\log(2r),
 \qquad
 d=\max\left\{3,\left\lceil\frac{\ell_r}{2}\right\rceil\right\},
 \qquad
 h=\log(2rd).
\]
Then $d$ is an integer and
$
 3\le d\le\max\{3,\ell_r\},
$
so $d$ is admissible in \Cref{thm:parameterised-main}. Indeed, if
$\ell_r<3$, then $d=3$. If $\ell_r\ge3$, then
$\lceil\ell_r/2\rceil\le\ell_r$.

Let
\[
 C_*=\frac3{\log4}.
\]
Since $\ell_r\ge\log4$, we have
\[
 d\le C_*\ell_r,
 \qquad
 d-1\ge\frac{\ell_r}{3}.
\]
For the second inequality, note that $d=3$ when $\ell_r\le6$, while for
$\ell_r>6$,
\[
 d-1\ge\frac{\ell_r}{2}-1\ge\frac{\ell_r}{3}.
\]
Moreover, $\log d\le\ell_r$, and hence
\[
 \ell_r\le h=\ell_r+\log d\le2\ell_r,
 \qquad
 \log h\le\ell_r.
\]
The last inequality follows from $2\ell_r\le e^{\ell_r}$ for
$\ell_r\ge\log4$.

Write
\[
 D_{r,d}=
 r^{2d/(d-1)}d^{4d/(d-1)}h^{4/(d-1)}.
\]
Using
\[
 \frac{2d}{d-1}=2+\frac2{d-1},
 \qquad
 \frac{4d}{d-1}=4+\frac4{d-1},
\]
we obtain
\begin{align*}
 D_{r,d}
 &=r^2d^4
 \exp\!\left(
 \frac{2\log r+4\log d+4\log h}{d-1}
 \right)\le e^{30}C_*^4r^2\ell_r^4.
\end{align*}

Next, set
\[
 Q_{r,d}=
 r^{(2d^2-1)/(d-1)^2}
 d^{2d(2d-1)/(d-1)^2}
 h^{2d^2/(d-1)^2}.
\]
The exact identities
\[
 \frac{2d^2-1}{(d-1)^2}
 =2+\frac{4d-3}{(d-1)^2},
 \qquad
 \frac{2d(2d-1)}{(d-1)^2}
 =4+\frac{6d-4}{(d-1)^2},
\]
and
\[
 \frac{2d^2}{(d-1)^2}
 =2+\frac{4d-2}{(d-1)^2}
\]
give
\[
 Q_{r,d}=r^2d^4h^2e^E,
\]
where
\[
 E=
 \frac{(4d-3)\log r+(6d-4)\log d+(4d-2)\log h}
 {(d-1)^2}.
\]
For $d\ge3$,
\[
 4d-3\le6(d-1),
 \qquad
 6d-4\le8(d-1),
 \qquad
 4d-2\le5(d-1).
\]
Therefore
\[
 E\le\frac{19\ell_r}{d-1}\le57,
\]
and consequently
\[
 Q_{r,d}\le4e^{57}C_*^4r^2\ell_r^6.
\]

Let $c_{\mathrm{par}},K_{\mathrm{par}}>0$ be the constants in
\Cref{thm:parameterised-main}. If
$
 k\ge K_{\mathrm{par}}(4e^{57}C_*^4)r^2\ell_r^6,
$
then the threshold in that theorem is satisfied. It follows that
\begin{align*}
 R_r(k)
 \le\exp\!\left(-c_{\mathrm{par}}\frac{k}{D_{r,d}}\right)r^{rk}\le\exp\!\left(
 -\frac{c_{\mathrm{par}}}{e^{30}C_*^4}
 \frac{k}{r^2\ell_r^4}
 \right)r^{rk}.
\end{align*}
Renaming the two absolute constants proves \Cref{thm:main}.
\end{proof}

\begin{thebibliography}{99}

\bibitem{BalisterEtAl2026}
P. Balister, B. Bollob\'{a}s, M. Campos, S. Griffiths, E. Hurley,
R. Morris, J. Sahasrabudhe, M. Tiba,
Upper bounds for multicolour Ramsey numbers,
\emph{J. Amer. Math. Soc.} 39(3) (2026), 765--780.

\bibitem{CamposEtAl2026}
M. Campos, S. Griffiths, R. Morris, J. Sahasrabudhe,
An exponential improvement for diagonal Ramsey,
\emph{Ann. of Math.} 203(3) (2026), 869--932.

\bibitem{Conlon2009}
D. Conlon,
A new upper bound for diagonal Ramsey numbers,
\emph{Ann. of Math.} 170(2) (2009), 941--960.

\bibitem{ConlonFoxSudakov2015}
D. Conlon, J. Fox, B. Sudakov,
Recent developments in graph Ramsey theory,
in \emph{Surveys in Combinatorics 2015},
A. Czumaj, A. Georgakopoulos, D. Kr\'{a}l, V. Lozin, O. Pikhurko, eds.,
London Math. Soc. Lecture Note Ser. 424,
Cambridge Univ. Press, Cambridge, 2015, 49--118.

\bibitem{Erdos1947}
P. Erd\H{o}s,
Some remarks on the theory of graphs,
\emph{Bull. Amer. Math. Soc.} 53(4) (1947), 292--294.

\bibitem{ErdosSzekeres1935}
P. Erd\H{o}s, G. Szekeres,
A combinatorial problem in geometry,
\emph{Compositio Math.} 2 (1935), 463--470.

\bibitem{GuptaEtAl2024}
P. Gupta, N. Ndiaye, S. Norin, L. Wei,
Optimizing the CGMS upper bound on Ramsey numbers,
arXiv:2407.19026 [math.CO], 2024.

\bibitem{NarangTang2026}
I. Narang, Y. Tang,
Schrijver number quasi-tensorization and multicolor Ramsey bounds via
robust OR polynomials,
arXiv:2607.25023 [math.CO], 2026.

\bibitem{Ramsey1930}
F. P. Ramsey,
On a problem of formal logic,
\emph{Proc. London Math. Soc.} (2) 30(1) (1930), 264--286.

\bibitem{Sah2023}
A. Sah,
Diagonal Ramsey via effective quasirandomness,
\emph{Duke Math. J.} 172(3) (2023), 545--567.

\bibitem{Thomason1988}
A. Thomason,
An upper bound for some Ramsey numbers,
\emph{J. Graph Theory} 12(4) (1988), 509--517.

\bibitem{Wigderson2025}
Y. Wigderson,
Upper bounds on diagonal Ramsey numbers
(after Campos, Griffiths, Morris, and Sahasrabudhe),
\emph{Ast\'{e}risque} 462 (2025), Exp. No. 1230, 85--138.

\end{thebibliography}
\end{document}
